% Standalone manuscript: no external inputs, bibliography, or figure files.
% The two Python certificates are embedded verbatim in the appendices.
\documentclass[11pt,a4paper]{article}
\usepackage[T1]{fontenc}
\usepackage{lmodern,microtype,amsmath,amssymb,amsthm,mathtools}
\usepackage[margin=24mm,headheight=15pt]{geometry}
\usepackage{booktabs,array,xcolor,fancyhdr,titlesec,xurl,listings}
\usepackage[colorlinks=true,linkcolor=teal!65!black,citecolor=teal!65!black,urlcolor=teal!65!black,bookmarksnumbered=true]{hyperref}
\definecolor{ink}{HTML}{15364A}
\definecolor{accent}{HTML}{007C78}
\definecolor{muted}{HTML}{566875}
\definecolor{pale}{HTML}{EEF6F5}
\titleformat{\section}{\Large\sffamily\bfseries\color{ink}}{\thesection}{.55em}{}
\titleformat{\subsection}{\large\sffamily\bfseries\color{accent}}{\thesubsection}{.55em}{}
\newtheorem{theorem}{Theorem}[section]
\newtheorem{lemma}[theorem]{Lemma}
\newtheorem{proposition}[theorem]{Proposition}
\newtheorem{corollary}[theorem]{Corollary}
\theoremstyle{definition}\newtheorem{definition}[theorem]{Definition}
\theoremstyle{remark}\newtheorem{remark}[theorem]{Remark}
\newcommand{\R}{\mathbb R}\newcommand{\Z}{\mathbb Z}
\newcommand{\C}{\mathbb C}\newcommand{\N}{\mathbb N}
\newcommand{\one}{\mathbf1}
\DeclareMathOperator{\dist}{dist}
\numberwithin{equation}{section}
\allowdisplaybreaks[2]
\pagestyle{fancy}\fancyhf{}
\fancyhead[L]{\small\sffamily Prime phases, polygons, and Nyman approximation}
\fancyhead[R]{\small\sffamily Further transfers}
\fancyfoot[L]{\footnotesize\sffamily Prepared for Keio | 20 November 2026}
\fancyfoot[R]{\small\thepage}
\renewcommand{\headrulewidth}{.3pt}
\setlength{\parindent}{1.1em}\setlength{\parskip}{3pt}
\setlength{\emergencystretch}{2em}\raggedbottom
\setcounter{tocdepth}{1}
\lstset{language=Python,basicstyle=\ttfamily\footnotesize,
  keywordstyle=\color{accent}\bfseries,commentstyle=\color{muted},
  stringstyle=\color{ink},showstringspaces=false,columns=fullflexible,
  keepspaces=true,breaklines=true,breakatwhitespace=false,
  frame=single,rulecolor=\color{accent!25},framesep=5pt,
  numbers=left,numberstyle=\tiny\color{muted},numbersep=8pt,
  xleftmargin=12pt,aboveskip=8pt,belowskip=8pt}
\hypersetup{pdftitle={Prime Phase Rigidity, Polygon Positivity, and an Exact Nyman Cone Projection},pdfauthor={Shaitetsu; research manuscript prepared for Keio},pdfsubject={Complete further closure attempts from VMRC 200; no proof of RH claimed}}
\begin{document}
\begin{titlepage}
\noindent{\sffamily\small\bfseries\color{accent} GEOMETRY OF NUMBERS \;|\; ANALYTIC NUMBER THEORY}\par
\vspace{9mm}
\noindent{\sffamily\bfseries\fontsize{25}{29}\selectfont\color{ink}\raggedright
Prime Phase Rigidity,\\
Polygon Positivity,\\
and an Exact Nyman\\
Cone Projection\par}
\vspace{5mm}
\noindent{\large\sffamily Further closure attempts from lattice covering flows}\par
\vspace{7mm}
\noindent Shaitetsu\quad$\cdot$\quad Research manuscript prepared for Keio\par
\noindent 20 November 2026\par
\vspace{6mm}
\begingroup\small
\noindent\textbf{Abstract.}
We continue three exact transfers from the positive fractional-part
flow of VMRC [200] to questions equivalent to the Riemann hypothesis.
For the original reciprocal-prime phases, we certify a common arc
of length less than $0.49$ and solve the integer-lift minimization
in every dimension: the closest quotient-lattice vector is uniquely
zero. Both cubic losses become explicit weighted pair sums.
The upper bound retaining all jump-midpoint corrections still
diverges unconditionally on the actual prime data.

For Li curvature, a verified zero height $H$ yields positive
Toeplitz matrices in every dimension after damping by
$\exp(-|n|\operatorname{arsinh}(1/(2H^2)))$.
Two populated bands of verified zeros also give strict positivity
of the original undamped Fejer polynomials for every degree
$M\le H/10$. Published inputs imply this for $M\le3\cdot10^{11}$.
Radial projection, finite angular averaging, and minimal damping
are examined with their exact limitations.

We determine the full nonnegative Nyman cone projection using
exactly the eight indices $\{2,3,4,5,7,11,13,17\}$.
Its squared distance is enclosed by
$0.0273089278984818<\mathcal E_+<0.0273089278984835$.
A finite integer-interval certificate and a monotone analytic tail
certify all infinitely many excluded generators.
Signed approximants satisfy explicit negative-mass laws.
A positive event lift is constructed, and a relaxed zero-error
limit is shown to violate the original finite divisor support.
\par\endgroup
\vfill
\noindent\fcolorbox{accent!35}{pale}{\begin{minipage}{.94\textwidth}
\textbf{What is closed}\par\smallskip
The prime closest-point choice in every dimension; a large finite
range of original polygon positivity; every degree after explicit
damping; and the entire nonnegative Nyman cone optimization.\par\medskip
\textbf{What remains}\par\smallskip
Uniform prime variance, positivity at all original polygon orders,
and vanishing signed Nyman error remain unproved.
No proof or disproof of the Riemann hypothesis is claimed.
\end{minipage}}
\end{titlepage}

\tableofcontents
\medskip
\noindent\textbf{Reading guide.}
Sections~\ref{r5p:section}--\ref{r5p:losssection} treat the actual
prime phases and the cubic losses.
Section~\ref{r5l:section} establishes both positive regularization
and a finite undamped range.
Section~\ref{r5g:section} closes the global nonnegative cone problem.
Section~\ref{r5n:sec} quantifies signed cancellation and analyzes
positive lifts.
The appendices include the analytic prime criterion and complete
integer-arithmetic verifiers. This single source file compiles
without external figures, inputs, or bibliography files.
\clearpage

\section{The three transfers and the scope of this continuation}
\label{r5:intro}

The objective is to push the preceding lattice-to-zeta transfers
past a change of notation and prove concrete estimates for the
actual arithmetic objects. This manuscript contains the complete
proofs and certificates of the present research batch. It is a
separate continuation of \cite{TransferZeta}, not a reprint of
every earlier manuscript.
The corpus sources [149] and [200] were checked at revision
\texttt{ac1492a50784be919fe0b5e8f19f96f8d221a22e}
\cite{PrimeCircles,VMRC200}.

The conclusions concern specified geometric and approximation
problems. They do not assert historical priority or resolve the
unrestricted arithmetic statements merely by transporting them.
Classical analytic inputs, including the strong
B\'aez-Duarte theorem, positive-real moment criteria, verified
zero height, and explicit zero counting, are cited where used.

\subsection{The positive-flow mechanism of corpus [200]}

For positive integers $r_j$, positive real weights $c_j$, and real
phases $t_j$, set
\[
 H_t(s)=\sum_jc_j\{t_j+r_js\},\qquad C=\sum_jr_jc_j>0.
\]
At each event $e$, let $J_e>0$ be the total downward jump and
$H_t(e)$ its post-event value. Then
\[
 \int_0^1H_t(s)\,ds=\frac12\sum_jc_j,\qquad
 \sum_eJ_e=C.
\]
Balancing the continuous derivative and jumps of $H_t^2$ gives
\[
 0=2C\int_0^1H_t(s)\,ds
       -2\sum_eJ_eH_t(e)-\sum_eJ_e^2.
\]
The positivity of every event weight therefore proves the central
inequality of \cite{VMRC200}:
\begin{equation}\label{r5:flow}
 \min_sH_t(s)\le
 \frac12\sum_jc_j-\frac{\sum_jr_jc_j^2}{2C}.
\end{equation}
For coincident events, expanding $J_e^2$ only adds nonnegative
cross terms. The positive slope and positive averaging weights
are the features that must be tracked in every proposed transfer.

\subsection{Prime circles and their actual lattice points}

Let $p_n$ be the $n$th prime. The actual data of \cite{PrimeCircles}
are
\[
 \theta_n=\sum_{p\le p_n}\frac1p-\log\log p_n,\quad
 w_n=p_{n+1}-p_{n-1}\quad(n\ge2).
\]
With $W_N=\sum_{n=2}^Nw_n$, define
\[
 Q_N=\min_s\sum_{n=2}^Nw_n\|\theta_n+s\|_{\Z}^2,\qquad
 D_N=W_N-\left|\sum_{n=2}^Nw_ne^{2\pi i\theta_n}\right|.
\]
The elementary chord bounds
$16\|t\|_{\Z}^2\le4\sin^2(\pi t)\le4\pi^2\|t\|_{\Z}^2$
and minimization over a unit circle direction give
\begin{equation}\label{r5:primecomparison}
 8Q_N\le D_N\le2\pi^2Q_N.
\end{equation}
Appendix~\ref{r4a:appendix} proves the analytic equivalence
\[
 \mathrm{RH}\quad\Longleftrightarrow\quad
 \sup_ND_N<\infty
 \quad\Longleftrightarrow\quad \sup_NQ_N<\infty.
\]
The main new geometric step below computes the minimizer in $Q_N$
for every $N$. The required uniform bound on its value is separate.

\subsection{Zero moments and exact finite polygon tests}

Write $\xi(s)=\frac12s(s-1)\pi^{-s/2}\Gamma(s/2)\zeta(s)$ and
let $\rho$ run over its zeros with multiplicity. Define
\[
 u_\rho=1-\rho^{-1},\qquad
 w_\rho=[\rho(1-\rho)]^{-1},\qquad
 a_n=\sum_\rho w_\rho u_\rho^n\quad(n\in\Z).
\]
These absolutely convergent sums are real and even for each fixed
$n$. They are the second differences of the Li coefficients.
In particular $a_0=2+\gamma_E-\log(4\pi)$.
For $M\ge1$, put $L_M=2M+1$ and
\begin{equation}\label{r5:polygontests}
 \eta_{M,j}=\frac1{L_M}
 \left[a_0+2\sum_{n=1}^M
  \left(1-\frac n{M+1}\right)a_n
       \cos\frac{2\pi nj}{L_M}\right],
 \quad 0\le j<L_M.
\end{equation}
The exact criterion of \cite{TransferZeta} is
\begin{equation}\label{r5:polygoncriterion}
 \mathrm{RH}\quad\Longleftrightarrow\quad
 \eta_{M,j}\ge0\text{ for every }M,j.
\end{equation}
Here is a brief independent verification. Under RH, $|u_\rho|=1$
and $w_\rho=(1/4+(\Im\rho)^2)^{-1}>0$; the $a_n$ are moments of
a positive circle measure, and its Fejer convolutions are
nonnegative. Conversely, if all weights in
\eqref{r5:polygontests} are nonnegative, they define discrete
positive measures of mass $a_0$. For each $|n|\le M$, their
$n$th moment is exactly $(1-|n|/(M+1))a_n$, by the finite
Fourier identity at $2M+1$ nodes. Letting $M\to\infty$ in every
fixed finite quadratic form makes $(a_{k-j})$ positive
semidefinite, hence $|a_n|\le a_0$.

If RH failed, zero symmetry would give a node with $|u_\rho|>1$.
The germ
\[
 A(z)=\sum_{n\ge0}a_nz^n
     =\sum_\rho\frac{w_\rho}{1-u_\rho z}
\]
would have an uncancelled pole at $1/u_\rho$ inside the unit
disk. The Cayley map is injective, so only the positive integer
multiplicity of that zero can share this pole, and its coefficient
is nonzero. The tail is locally normally convergent away from
poles, since $u_\rho\to1$ and $\sum_\rho|w_\rho|<\infty$.
This contradicts bounded Taylor coefficients. The moment and
positive-real interpretation is classical \cite{Lagarias}.
The finite-degree theorem below proves a substantial bounded range
of \eqref{r5:polygoncriterion}, not its universal quantifier.

\subsection{The signed span and the nonnegative cone}

For $n\ge2$, set
\[
 f_n(x)=\{1/(nx)\}-\{1/x\}/n,\qquad 0<x<1.
\]
The strong B\'aez-Duarte theorem gives
\begin{equation}\label{r5:nymancriterion}
 \mathrm{RH}\quad\Longleftrightarrow\quad
 \delta_N:=\inf_{c_2,\ldots,c_N\in\R}
       \left\|1-\sum_{n=2}^Nc_nf_n\right\|_{L^2(0,1)}
       \longrightarrow0
\end{equation}
\cite{BaezDuarte}. To check the normalization, let
$e_n(x)=\{1/(nx)\}$ on $(0,\infty)$.
Each $f_n=e_n-e_1/n$ vanishes on $(1,\infty)$.
Conversely, if $h=\sum_{n=1}^Nd_ne_n$ approximates
$1_{(0,1)}$ there, its tail is $A/x$, where
$A=\sum d_n/n$, so $|A|\le\|1_{(0,1)}-h\|_2$.
Replacing $h$ by $h-Ae_1=\sum_{n=2}^Nd_nf_n$ changes the
error by at most $|A|\|e_1\|_2$. This proves the stated
equivalence of the two approximation normalizations.

Restricting every $c_n$ to be nonnegative gives a smaller problem.
We determine its exact global optimum in
Theorem~\ref{r5g:global}. Its positive minimum is compatible with
\eqref{r5:nymancriterion}, which permits signed coefficients.
The later mass estimates quantify how much negative coefficient
mass any successful signed sequence must develop.


\section{The actual prime point: a unique lift in every dimension}
\label{r5p:section}

We retain exactly the phases and weights of \cite{PrimeCircles}:
\begin{equation}\label{r5p:data}
 \theta_n=\sum_{p\leq p_n}\frac1p-\log\log p_n,
 \qquad w_n=p_{n+1}-p_{n-1},\qquad
 W_N=\sum_{n=2}^Nw_n=p_N+p_{N+1}-5.
\end{equation}
Let $\bar\theta_N=W_N^{-1}\sum_{n=2}^Nw_n\theta_n$.
The quotient-lattice quantity from the preceding investigation is
\[
 Q_N=\min_s\sum_{n=2}^Nw_n\|\theta_n+s\|_{\Z}^2
     =\operatorname{dist}_w(\pi_w\theta,\Lambda_w)^2,
 \quad
 \pi_wx=x-\frac{\sum_nw_nx_n}{W_N}\mathbf1,
 \quad \Lambda_w=\pi_w\Z^{N-1}.
\]
The weighted inner product is
$\langle x,y\rangle_w=\sum_nw_nx_ny_n$.
We first resolve its integer-lift and common-rotation minimizations
for every original prime prefix.  We then compute both cubic losses
explicitly and show that the entire midpoint-corrected upper bound
still diverges on the actual prime data.

\subsection{A certified arc containing all the prime phases}

\begin{theorem}[Global prime-phase arc]\label{r5p:arc}
For every prime $p\geq3$,
\begin{equation}\label{r5p:arcinterval}
 \frac14<\sum_{q\leq p}\frac1q-\log\log p<\frac{37}{50}.
\end{equation}
In particular, all real pair differences of the phases in
\eqref{r5p:data} are strictly smaller than $49/100$ in absolute
value.
\end{theorem}

\begin{proof}
Write $\theta(x)=\sum_{p\leq x}p^{-1}-\log\log x$ and let
$B_{\mathbb P}=\lim_{x\to\infty}\theta(x)$.
Rosser--Schoenfeld \cite[Thm.~5, eqs.~(3.17)--(3.18)]{RosserSchoenfeld}
give the unconditional estimate
\begin{equation}\label{r5p:RS}
 |\theta(x)-B_{\mathbb P}|<\frac1{2\log^2x}
 \qquad(x\geq286).
\end{equation}
Set $T=10000$.  Eliminating $B_{\mathbb P}$ between
\eqref{r5p:RS} at $T$ and at any $x\geq T$ gives
\begin{equation}\label{r5p:tailarc}
 \theta(T)-\frac1{\log^2T}
 <\theta(x)<
 \theta(T)+\frac1{\log^2T}.
\end{equation}
Thus a finite calculation suffices, with no numerical input value
of $B_{\mathbb P}$.

The exact-integer certificate described below gives the following
enclosures.  Displayed decimal endpoints have been rounded outward.
There are $1229$ primes up to $T$, of which $1228$ are at least $3$.
\begin{center}
\begin{tabular}{@{}lll@{}}
\hline
Quantity & Lower bound & Upper bound\\
\hline
Finite minimum, at $p=8597$ & $0.262115904496$ & $0.262115904497$\\
Finite maximum, at $p=3$ & $0.739285505716$ & $0.739285505717$\\
$\theta(10000)$ & $0.262733140865$ & $0.262733140866$\\
All $\theta(x)$, $x\geq10000$ & $0.250944909802$ & $0.274521371929$\\
\hline
\end{tabular}
\end{center}
Both the finite enclosures and \eqref{r5p:tailarc} lie strictly in
$(1/4,37/50)$, proving the theorem.

Here are the details that make the finite step rigorous.  All
interval endpoints are integers divided by $D=10^{18}$.
For a rational $r\in[1,2]$, put $z=(r-1)/(r+1)\in[0,1/3]$ and use
\begin{equation}\label{r5p:logcertificate}
 \log r=2\sum_{j=0}^{23}\frac{z^{2j+1}}{2j+1}+R,
 \qquad
 0\leq R\leq\frac9{4\cdot49\cdot3^{49}}<\frac1D.
\end{equation}
Round each series term downward by exact integer division.  The
upper endpoint adds $48$ rounding units and one remainder unit.
Reduction $r=2^km$, $1\leq m<2$, handles every positive rational
argument, with the same enclosure for $\log2$.
For $\log\log p$, apply monotonicity and these logarithm bounds
to the two rational endpoints enclosing $\log p$.
Enclose $\sum_{p\leq x}1/p$ by adding integer floors of $D/p$,
with one upper rounding unit per prime.  Finally bound
$1/\log^2T$ from above using the lower enclosure of $\log T$.
An Eratosthenes sieve and these formulas require only integer
arithmetic; the complete program in Appendix~\ref{r5:prime-code}
records all constants and assertions. Formula \eqref{r5p:logcertificate}
provides the explicit error bound for every transcendental operation.
\end{proof}

\subsection{A sharp short-arc rigidity inequality}

For arbitrary positive weights $w_1,\ldots,w_m$, set
\[
 \mathcal V_w(t)=\sum_iw_i(t_i-\bar t)^2
       =\frac1W\sum_{i<j}w_iw_j(t_i-t_j)^2,
 \qquad \bar t=\frac{\sum_iw_it_i}{W}.
\]
The central Voronoi cell is denoted by
\[
 \operatorname{Vor}_w(\Lambda_w)
 =\{x\in H_w:\|x\|_w\leq\|x-\lambda\|_w
                  \text{ for every }\lambda\in\Lambda_w\},
 \quad H_w=\{x:\textstyle\sum_iw_ix_i=0\}.
\]

\begin{theorem}[Uniform integer-lift rigidity]\label{r5p:rigidity}
Suppose the real numbers $t_i$ have span at most $L<1/2$.
For every $z\in\Z^m$,
\begin{equation}\label{r5p:rigidityineq}
 \mathcal V_w(t-z)-\mathcal V_w(t)
       \geq(1-2L)\mathcal V_w(z).
\end{equation}
Consequently
\begin{equation}\label{r5p:exactlift}
 \min_s\sum_iw_i\|t_i+s\|_{\Z}^2=\mathcal V_w(t),
 \qquad s=-\bar t\pmod1
\end{equation}
is the unique optimal common shift.  Every optimal integer lift is
constant, and the unique closest lattice vector to $\pi_wt$ is $0$.
Moreover,
\begin{equation}\label{r5p:scaledVor}
 \pi_wt\in2L\operatorname{Vor}_w(\Lambda_w).
\end{equation}
The factor $1-2L$ in \eqref{r5p:rigidityineq} is sharp.
\end{theorem}

\begin{proof}
For $d=t_i-t_j$ and $k=z_i-z_j\in\Z$,
\[
 (d-k)^2-d^2=k^2-2dk
       \geq k^2-2L|k|\geq(1-2L)k^2.
\]
Summing the pair formula proves \eqref{r5p:rigidityineq}.
Joint minimization over $s$ and $z$ first minimizes the real
quadratic in $s$, giving $\mathcal V_w(t-z)$.
Since $1-2L>0$, its minimum occurs exactly when $z$ is constant.
For those lifts the shift is $-\bar t$ modulo $1$, proving
\eqref{r5p:exactlift} and uniqueness.

For $\lambda=\pi_wz$, inequality \eqref{r5p:rigidityineq} also gives
\[
 2\langle\pi_wt,\lambda\rangle_w\leq2L\|\lambda\|_w^2.
\]
The half-space description of the central Voronoi cell now gives
\eqref{r5p:scaledVor}.  The case $L=0$ is immediate.
For sharpness, take two points $t=(0,L)$ and $z=(0,1)$;
the pair calculation is an equality for arbitrary positive weights.
\end{proof}

\begin{lemma}[Shortest nonzero quotient vector]\label{r5p:shortest}
If $m\geq2$ and $w_{\min}=\min_iw_i$, then
\begin{equation}\label{r5p:shortestformula}
 \min_{0\ne\lambda\in\Lambda_w}\|\lambda\|_w^2
       =\frac{w_{\min}(W-w_{\min})}{W}.
\end{equation}
\end{lemma}

\begin{proof}
For nonconstant $z\in\Z^m$, choose a nontrivial level cut with
weight $A$.  The integer differences across the cut have absolute
value at least $1$, so
\[
 \mathcal V_w(z)\geq\frac{A(W-A)}W
       \geq\frac{w_{\min}(W-w_{\min})}W.
\]
Both sides of the cut have weight at least $w_{\min}$.
Equality is attained by $z=e_i$ at an index of minimum weight.
\end{proof}

\begin{corollary}[The original prime closest point is known exactly]
\label{r5p:primeclosest}
For every $N\geq2$,
\begin{equation}\label{r5p:primevariance}
 \boxed{\displaystyle
 Q_N=\sum_{n=2}^Nw_n(\theta_n-\bar\theta_N)^2
     =\sum_{n=2}^Nw_n\theta_n^2
          -\frac{(\sum_{n=2}^Nw_n\theta_n)^2}{W_N}.}
\end{equation}
The optimal shift is $-\bar\theta_N$ modulo $1$, the unique closest
lattice vector is $0$, and
$\pi_w\theta\in\operatorname{int}((49/50)
\operatorname{Vor}_w(\Lambda_w))$ in the relative ambient space.
For $N\geq3$, every nonconstant integer lift satisfies
\begin{equation}\label{r5p:primegap}
 \mathcal V_w(\theta-z)-Q_N
       \geq\frac{3(W_N-3)}{50W_N}\geq\frac6{175}.
\end{equation}
The weighted Euclidean distance of $\pi_w\theta$ from every
Voronoi boundary hyperplane is at least
\begin{equation}\label{r5p:clearance}
 \frac1{100}\sqrt{\frac{3(W_N-3)}{W_N}}
      \geq\frac1{100}\sqrt{\frac{12}{7}}.
\end{equation}
\end{corollary}

\begin{proof}
Apply Theorem~\ref{r5p:rigidity} with $L=49/100$ using
Theorem~\ref{r5p:arc}.  The actual span is strictly smaller than
$L$, giving the stated relative interior inclusion.
For the prime weights, $w_2=3$ and $w_n\geq4$ for $n\geq3$.
Also $W_N\geq7$ for $N\geq3$.
Inequality \eqref{r5p:primegap} follows from
\eqref{r5p:rigidityineq} and Lemma~\ref{r5p:shortest}.
For every nonzero lattice vector $\lambda$, the distance to its
bisector is
\[
 \frac{\|\lambda\|_w^2-2\langle\pi_w\theta,\lambda\rangle_w}
      {2\|\lambda\|_w}
 \geq\frac{\|\lambda\|_w}{100}.
\]
Another use of Lemma~\ref{r5p:shortest} proves the clearance bound.
\end{proof}

This resolves the lift selection for every prefix, without a
nonresonance hypothesis about $B_{\mathbb P}$.  Weighted Cesaro
convergence gives $\bar\theta_N\to B_{\mathbb P}$, so the optimal
rotations tend to $-B_{\mathbb P}$ modulo $1$.
The remaining unknown is the size of the weighted norm in
\eqref{r5p:primevariance}.  A fixed fractional depth inside a
growing Voronoi cell does not bound that norm uniformly.

\section{The exact cubic losses and an arithmetic obstruction}
\label{r5p:losssection}

Define the weighted pairwise absolute spread
\begin{equation}\label{r5p:Gini}
 \mathcal A_N=\frac1{W_N}
       \sum_{2\leq i<j\leq N}w_iw_j|\theta_i-\theta_j|.
\end{equation}
For this section fix $N$, suppress its index on the weights, and
put $x_i=\theta_i-\bar\theta_N$.  Use the periodic flow
\[
 q_N(s)=\sum_iw_i\|x_i+s\|_{\Z}^2,
 \quad F_N(s)=q_N'(s)/2,
 \quad G_N(s)=q_N(s)-F_N(s)^2/W_N
\]
away from its crossings.  We have $q_N(0)=Q_N$ and $F_N(0)=0$.
Let
\[
 R_N^2=\frac{W_N}{12}-\frac{\sum_iw_i^3}{12W_N^2}.
\]
At the jump of size $w_i$, write $y_i$ for the midpoint of the
pre-jump and post-jump values of $F_N$.  Simultaneous jumps are
resolved sequentially using formal intermediate values, as in the
cubic jump identity.

\begin{theorem}[Both losses on the actual prime arc]
\label{r5p:losses}
The midpoint-jump loss and the branch loss satisfy exactly
\begin{align}
 L_{\rm jump}
   :=\frac1{W_N^2}\sum_iw_iy_i^2
   &=R_N^2+Q_N-\mathcal A_N,\label{r5p:jumploss}\\
 L_{\rm branch}
   :=\int_0^1(G_N(s)-Q_N)\,ds
   &=\mathcal A_N-2Q_N.\label{r5p:branchloss}
\end{align}
In particular, the upper bound retaining every jump-midpoint
correction, and discarding only the branch loss, is
\begin{equation}\label{r5p:correctedupper}
 Q_N\leq\mathcal U_N:=R_N^2-L_{\rm jump}
                =\mathcal A_N-Q_N.
\end{equation}
\end{theorem}

\begin{proof}
Sort $x_1\geq\cdots\geq x_m$, and put
$B_{i-1}=\sum_{j<i}w_j$ and
$c_i=(B_{i-1}+w_i/2)/W_N$.
Every $x_i$ has absolute value less than $1/2$.
Its crossing is at $s_i=1/2-x_i\in(0,1)$, in the chosen order,
and therefore
\[
 \frac{y_i}{W_N}=\frac12-c_i-x_i.
\]
Equal $x_i$ cause simultaneous crossings; the formula still gives
their successive formal midpoints in any ordering of the ties.

Consider the partition of $[0,1]$ into consecutive lengths
$w_i/W_N$, with midpoints $c_i$.  Exact midpoint quadrature for
$(u-1/2)^2$ gives
\[
 \sum_iw_i(1/2-c_i)^2
       =\frac{W_N}{12}-\frac{\sum_iw_i^3}{12W_N^2}=R_N^2.
\]
Furthermore, $\sum_iw_ix_i^2=Q_N$ and direct summation gives
\[
 2\sum_iw_ix_i(1/2-c_i)
  =\frac1{W_N}\sum_{i<j}w_iw_j(x_i-x_j)=\mathcal A_N.
\]
Expanding $\sum_iw_i(1/2-c_i-x_i)^2$ proves
\eqref{r5p:jumploss}.

There is a direct pair proof of \eqref{r5p:branchloss}.
For a pair $x_i\geq x_j$, set $d=x_i-x_j<1/2$.
Its rounded integer difference is $1$ exactly between its two
crossings, on an interval of length $d$, and is $0$ elsewhere.
On that interval the increase of its squared lifted difference
is $(d-1)^2-d^2=1-2d$.
The weighted pair formula for $G_N$, followed by integration,
therefore gives
\[
 \int_0^1(G_N(s)-Q_N)\,ds
   =\frac1{W_N}\sum_{i<j}w_iw_j(d-2d^2)
   =\mathcal A_N-2Q_N.
\]
Ties have $d=0$ and contribute zero.  Combining the two exact
losses gives $Q_N=R_N^2-L_{\rm jump}-L_{\rm branch}$,
and hence \eqref{r5p:correctedupper}.
\end{proof}

The next result concerns the actual prime data, not an auxiliary
model sequence.  It shows that even the corrected upper bound
\eqref{r5p:correctedupper} cannot itself be the desired constant
bound.

\begin{theorem}[Unconditional divergence of the corrected cubic bound]
\label{r5p:divergence}
For the original phases and weights,
\begin{equation}\label{r5p:divergenceformula}
 \mathcal A_N\longrightarrow\infty,
 \qquad
 \mathcal U_N\geq\frac{51}{100}\mathcal A_N
             \longrightarrow\infty.
\end{equation}
Neither conclusion assumes RH or its negation.
\end{theorem}

\begin{proof}
We first prove that
\begin{equation}\label{r5p:sampledL1}
 \sum_{n=2}^{\infty}w_n|E(p_n)|=\infty,
 \qquad E(x)=\theta(x)-B_{\mathbb P}.
\end{equation}
Suppose the sum were finite, and write $e_n=E(p_n)$,
$p=p_n$, $q=p_{n+1}$, and $h_n=q-p$.
On $[p,q)$ the function $E$ is a linear interpolation in
$\log\log x$ between $e_n$ and $e_{n+1}-1/q$, whence
\[
 |E(x)|\leq|e_n|+|e_{n+1}|+1/q.
\]
For every $\sigma>0$ the endpoint terms in
$\int_2^\infty|E(x)|x^{-\sigma}\,dx$ are summable by the
assumption and the identity $w_n=h_{n-1}+h_n$.
Crucially, although the unweighted $L^1$ jump correction is not
summable, its Mellin-weighted version is.  Bertrand's postulate gives
\begin{equation}\label{r5p:weightedcorrection}
 \sum_{n\geq1}\frac{h_n}{p_{n+1}p_n^\sigma}
 \leq2^\sigma\sum_{n\geq1}\frac{h_n}{p_{n+1}^{1+\sigma}}
 \leq2^\sigma\int_2^\infty x^{-1-\sigma}\,dx
 =\frac1\sigma.
\end{equation}
Thus $F(s)=\int_2^\infty E(x)x^{-s}\,dx$ is holomorphic for
$\Re s>0$.  For differentiation, any power of $\log x$ can be
absorbed by taking a slightly smaller positive real exponent.

The prime-square obstruction can now be applied without an
unweighted $L^1$ sampling claim.  Put $c=\log\log2+B_{\mathbb P}$.
For $\Re s>1$, differentiation of the prime-zeta identity gives
\[
 P'(s)=H(s):=F(s)+(s-1)F'(s)
           -c(\log2)2^{1-s}-\frac{2^{1-s}}{s-1},
 \qquad P(s)=\sum_p p^{-s}.
\]
Under the assumed sampled bound, $H$ is meromorphic for
$\Re s>0$ with its only pole at $1$, of residue $-1$.
The Euler product and meromorphic continuation give, for
$\Re s>1/2$,
\begin{equation}\label{r5p:squareobstruction}
 \frac{\zeta'(s)}{\zeta(s)}
       =H(s)+H(2s)+\sum_{k\geq3}P'(ks).
\end{equation}
The $k\geq3$ sum is holomorphic in a neighborhood of the real
point $1/2$, by absolute local uniform convergence.
As real $s\downarrow1/2$, $H(s)$ stays bounded but
\[
 H(2s)=-\frac1{2(s-1/2)}+O(1).
\]
This is impossible: $\zeta(1/2)\ne0$, since the alternating eta
series is positive for real $0<s<1$ and
$\zeta(s)=\eta(s)/(1-2^{1-s})<0$ there.
Thus $\zeta'/\zeta$ is bounded near that real point.
The contradiction proves \eqref{r5p:sampledL1}.
This is the prime-square Mellin mechanism of \cite{PrimeCircles},
with the weighted endpoint correction made explicit.

Next, for any bounded convergent real sequence $t_i\to B$ and
positive weights with $W_N\to\infty$, Jensen's inequality gives
\begin{equation}\label{r5p:GiniJensen}
 \sum_{i\leq N}w_i|t_i-\bar t_N|
 \leq\frac1{W_N}\sum_{i,j\leq N}w_iw_j|t_i-t_j|
 =2\mathcal A_N.
\end{equation}
If $\mathcal A_N$ were bounded along an infinite subsequence,
pass to that subsequence in every fixed finite prefix of
\eqref{r5p:GiniJensen}.  Weighted Cesaro convergence
$\bar t_N\to B$ would imply $\sum_iw_i|t_i-B|<\infty$.
Apply this with $t_i=\theta_i$, $B=B_{\mathbb P}$, and
\eqref{r5p:sampledL1}.  No bounded subsequence is possible, so
$\mathcal A_N\to\infty$.

Finally, Theorem~\ref{r5p:arc} gives, term by term,
\[
 Q_N=\frac1{W_N}\sum_{i<j}w_iw_j(\theta_i-\theta_j)^2
       \leq\frac{49}{100}\mathcal A_N.
\]
Equation \eqref{r5p:correctedupper} now proves
\eqref{r5p:divergenceformula}.
\end{proof}

Under RH, $Q_N$ is bounded, but both
$\mathcal U_N=\mathcal A_N-Q_N$ and
$L_{\rm branch}=\mathcal A_N-2Q_N$ diverge.
Their difference is exactly the bounded target.  The jump identity
has therefore isolated an unresolved cancellation rather than
provided its uniform estimate.

\subsection{Relation to the earlier corpus and the remaining estimate}

The real weighted variance in \eqref{r5p:primevariance}, its
divergent-regime comparison with the circle deficit, and its
comparison with the continuous prime-counting variance already
appear in \cite{PrimeCircles}.  The new conclusion established here
is its exact identification with the torus and quotient-lattice
minimum for \emph{every} $N$, together with a certified lift, a
unique closest vector, and quantitative interior bounds.
The prime-square obstruction is also an existing analytic mechanism;
the proof above uses its weighted discrete form to analyze the
specific cubic upper bound.

What remains unproved is
\begin{equation}\label{r5p:remaining}
 \sup_N\sum_{n=2}^Nw_n(\theta_n-\bar\theta_N)^2<\infty.
\end{equation}
The analytic bridge makes \eqref{r5p:remaining} equivalent to RH.
All integer and rotation optimization in this expression has now
been resolved.  The present results do not bound the remaining real
weighted variance, and no historical-priority claim is attached to
the auxiliary geometric formulas.

\section{Li curvature: positive regularization and a finite undamped range}
\label{r5l:section}

The ordinary Li criterion and the positive-real logarithmic-derivative
criterion are classical \cite{Lagarias}.  We retain the curvature
normalization and exact polygon transport of \cite{TransferZeta}.
The results below prove positivity for an explicitly damped sequence
at every degree, and then prove positivity for the original Fejer
polynomials through a large explicit finite degree.  The latter
argument uses two populated bands of verified critical-line zeros
to dominate all possible off-line contributions above the verified
height.  Neither result establishes the infinite undamped assertion.

\subsection{Normalization and the geometry of one zero pair}

For nontrivial zeros, counted with multiplicity, write
\begin{equation}\label{r5l:notation}
 u_\rho=1-\rho^{-1},\qquad
 w_\rho=[\rho(1-\rho)]^{-1},\qquad
 a_n=\sum_\rho w_\rho u_\rho^n\quad(n\in\Z).
\end{equation}
The sum converges absolutely for every fixed $n$.  The sequence is
real and even, and is the second difference of the ordinary Li
coefficients, with $\lambda_0=0$ and $\lambda_{-n}=\lambda_n$.
In particular,
\[
 a_0=2+\gamma_E-\log(4\pi)>0.
\]
The identities established in \cite{TransferZeta} are
\begin{align}
 A(z)&=\sum_{n\ge0}a_nz^n
     =\sum_\rho\frac{w_\rho}{1-u_\rho z},\label{r5l:Agerm}\\
 F(z)&=2A(z)-a_0
     =2\frac{\xi'}{\xi}\left(\frac1{1-z}\right).
       \label{r5l:Fgerm}
\end{align}
They initially hold near zero, and then by meromorphic continuation.
The conjugate-paired Hadamard product gives
$\xi'/\xi(s)=\lim_{T\to\infty}\sum_{|\Im\rho|\le T}(s-\rho)^{-1}$
without an additional constant.  Its real-part series converges
absolutely at every regular fixed $s$, which permits the horizontal
pairing used below.

We write $\mathrm{RH}(H)$ for the statement that all nontrivial zeros
with $|\Im\rho|\le H$ have real part $1/2$.  The published
Platt--Trudgian theorem supplies $H=3\cdot10^{12}$ by rigorous
interval arithmetic \cite{PlattTrudgian}.

\begin{lemma}[The negative disk of a horizontal pair]
\label{r5l:pairgeometry}
Let $\rho_\pm=1/2\pm\delta+i\gamma$, where $0<\delta<1/2$.
For $s=1/2+x+it$, $x>0$, away from the poles,
\begin{equation}\label{r5l:pairreal}
 \Re\left(\frac1{s-\rho_+}+\frac1{s-\rho_-}\right)
 =\frac{2x\{x^2+(t-\gamma)^2-\delta^2\}}
 {\{(x-\delta)^2+(t-\gamma)^2\}
  \{(x+\delta)^2+(t-\gamma)^2\}}.
\end{equation}
For $0<r<1$, the image of $|z|\le r$ under $s=1/(1-z)$
has center and radius
\begin{equation}\label{r5l:imagedisk}
 c_r=\frac1{1-r^2},\qquad v_r=\frac r{1-r^2},
 \qquad(c_r-1/2)^2-v_r^2=\frac14.
\end{equation}
This image is disjoint from the interior of the pair's negative
disk if and only if
\begin{equation}\label{r5l:disjoint}
 \gamma^2+\frac14-\delta^2\ge2v_r\delta.
\end{equation}
\end{lemma}
\begin{proof}
Adding the two real parts proves \eqref{r5l:pairreal}; its sign is
the sign of $x^2+(t-\gamma)^2-\delta^2$.  The image-circle formula
is an elementary computation.  External disjointness from the disk
with center $1/2+i\gamma$ and radius $\delta$ is
\[
 \sqrt{(c_r-1/2)^2+\gamma^2}\ge v_r+\delta.
\]
Squaring and using \eqref{r5l:imagedisk} gives
\eqref{r5l:disjoint}.  The image disk contains $s=1$, while the
negative disk does not, since $\delta<1/2$; hence containment cannot
provide another case of disjointness from its interior.
\end{proof}

\begin{proposition}[The exact Poisson threshold of one pair]
\label{r5l:pairthreshold}
For the Hermitian sequence
\[
 q_n=w_{\rho_+}u_{\rho_+}^n+w_{\rho_-}u_{\rho_-}^n,
\]
all finite Toeplitz matrices of $e^{-\tau|n|}q_n$ are positive
semidefinite if and only if
\begin{equation}\label{r5l:pairtau}
 \tau\ge\tau(\delta,\gamma):=
 \operatorname{arsinh}
       \frac{\delta}{\gamma^2+1/4-\delta^2}.
\end{equation}
A critical-line node has threshold zero.
\end{proposition}
\begin{proof}
For one zero,
\[
 \frac{w_\rho}{1-u_\rho z}
 =\frac1{1-\rho}+\frac1{s(z)-(1-\rho)}.
\]
The Caratheodory function $q_0+2\sum_{n\ge1}q_nz^n$ is therefore
twice the logarithmic-derivative sum for $1-\rho_+,1-\rho_-$,
plus a purely imaginary constant.  The reflected pair has center
$1/2-i\gamma$, so Lemma~\ref{r5l:pairgeometry} applies with the
same condition.  For $r=e^{-\tau}$,
\[
 v_r=\frac1{2\sinh\tau}.
\]
Thus \eqref{r5l:disjoint} is exactly \eqref{r5l:pairtau}.
The positive-real moment criterion proves sufficiency.  If the
condition fails, there is either a pole inside the disk or a
regular point with negative real part; both exclude a positive
circle moment representation.  At equality the possible tangency
is on the boundary, so positivity in the open disk, or closedness
of each finite PSD cone, includes the endpoint.
\end{proof}

This is an optimal threshold for a single horizontal pair.  It is
not asserted to be the optimal threshold of the full sum, where
other pairs can compensate for a negative contribution.

\subsection{An all-degree positive regularization}

\begin{theorem}[Positivity from a verified height]
\label{r5l:damped}
Assume $\mathrm{RH}(H)$ for $H>0$, and put
\begin{equation}\label{r5l:taudef}
 \tau_H=\operatorname{arsinh}\frac1{2H^2},\qquad
 r_H=e^{-\tau_H}
     =\frac{2H^2}{\sqrt{1+4H^4}+1}.
\end{equation}
For every $N\ge0$, the matrix
\begin{equation}\label{r5l:dampedmatrix}
 T_N^{(H)}=(r_H^{|k-j|}a_{k-j})_{j,k=0}^N
\end{equation}
is positive definite.  Its Fejer polynomials are strictly positive
at every angle.  Moreover, for every $n\ge0$,
\begin{equation}\label{r5l:allnbound}
 |a_n|\le a_0e^{n\tau_H}\le a_0e^{n/(2H^2)}.
\end{equation}
\end{theorem}
\begin{proof}
Every critical-line zero contributes a positive real part in
$\Re s>1/2$.  An off-line pair has $|\gamma|>H$ and
$0<\delta<1/2$.  For $r\le r_H$ we have $v_r\le H^2$, so
\[
 \gamma^2+1/4-\delta^2>H^2\ge2v_r\delta.
\]
Lemma~\ref{r5l:pairgeometry} applies strictly to each such pair.
The absolutely convergent real-part sum therefore shows that
$\Re F(r_Hz)>0$ on the closed unit disk.  No pole occurs there.
The function is holomorphic on a neighborhood of that disk:
strict disjointness and discreteness exclude nearby low zeros,
while zeros of sufficiently large height cannot approach its
compact image in the $s$-plane.

The boundary function $\Re F(r_He^{i\theta})$ is consequently a
strictly positive continuous density with respect to $d\theta/(2\pi)$, with moments
$r_H^{|n|}a_n$.  Integrating $|p|^2$ proves positive definiteness
of \eqref{r5l:dampedmatrix}.  Its Fejer convolutions are strictly
positive as well.  The principal two-by-two minors give
$|r_H^na_n|\le a_0$.  Finally $\operatorname{arsinh}x\le x$
proves the last inequality.
\end{proof}

Using the established height $H=3\cdot10^{12}$ gives
\begin{equation}\label{r5l:numerictau}
 0<\tau_H<\frac1{18\cdot10^{24}}
       =5.555\ldots\cdot10^{-26}.
\end{equation}
This number must not be rounded to zero.  For fixed $H$ the factor
$e^{n\tau_H}$ is unbounded as $n$ grows.

For the original arithmetic sequence define
\begin{equation}\label{r5l:fejerdef}
 F_M(\theta)=a_0+2\sum_{n=1}^M
    \left(1-\frac n{M+1}\right)a_n\cos(n\theta),
 \qquad
 \eta_{M,j}=\frac{F_M(2\pi j/(2M+1))}{2M+1}.
\end{equation}
\begin{corollary}[A bound for every original polygon]
\label{r5l:polygondeficit}
Under $\mathrm{RH}(H)$, for every $M\ge1$ and every vertex $j$,
\begin{equation}\label{r5l:polygondeficiteq}
 \eta_{M,j}\ge
 -\frac{a_0\tau_He^{M\tau_H}M(M+2)}{3(2M+1)}.
\end{equation}
\end{corollary}
\begin{proof}
The damped Fejer polynomial is nonnegative.  Its difference from
\eqref{r5l:fejerdef} has absolute value at most
\[
 2a_0\sum_{n=1}^M\left(1-\frac n{M+1}\right)
       (e^{n\tau_H}-1)
 \le2a_0\tau_He^{M\tau_H}
       \sum_{n=1}^M\left(1-\frac n{M+1}\right)n.
\]
The last sum is $M(M+2)/6$.  Division by $2M+1$ proves the bound.
\end{proof}

\subsection{A positive radial measure and its exact error}

For a zero $\rho=\beta+i\gamma$ in the critical strip,
\begin{equation}\label{r5l:positiveweights}
 \Re w_\rho=
 \frac{\gamma^2+\beta(1-\beta)}{|\rho(1-\rho)|^2}>0.
\end{equation}
Define
\begin{equation}\label{r5l:radialmeasure}
 \nu_{\rm rad}=\sum_\rho(\Re w_\rho)
                     \delta_{u_\rho/|u_\rho|},
 \qquad b_n=\int z^n\,d\nu_{\rm rad}(z).
\end{equation}
This is an unconditionally positive finite measure.  Its mass is
exactly $a_0$: imaginary parts cancel under the zero symmetries,
and taking real parts does not change the total mass.  Its moments
are real and even.  Reciprocal-conjugate nodes have the same radial
projection and conjugate weights, which explains this natural
choice of projected mass.

\begin{theorem}[A quantitative radial comparison]
\label{r5l:radialbound}
Assume $\mathrm{RH}(H)$, with $H\ge1/2$.  Put
\[
 W=\sum_{\Re\rho\ne1/2}\Re w_\rho\le a_0,
 \qquad
 d_n=\cosh(n\tau_H)-1+H^{-1}\sinh(n\tau_H).
\]
For all $n\ge0$,
\begin{equation}\label{r5l:radialerror}
 |a_n-b_n|\le Wd_n,
 \qquad |a_n|\le a_0+Wd_n.
\end{equation}
The original Toeplitz matrix $T_N=(a_{k-j})_{j,k=0}^N$ obeys
\begin{align}
 \lambda_{\min}(T_N)
 &\ge-2W\sum_{n=1}^Nd_n\label{r5l:radialmatrix}\\
 &\ge-W\cosh(N\tau_H)
 \left\{\frac{\tau_H^2N(N+1)(2N+1)}6
              +\frac{\tau_HN(N+1)}H\right\}.
       \label{r5l:radialmatrixexplicit}
\end{align}
\end{theorem}
\begin{proof}
In one horizontal off-line pair choose the member with $|u|>1$.
Write
\[
 u=Re^{i\theta},\qquad u^*=R^{-1}e^{i\theta},\qquad
 R=e^x>1,\qquad w=A+iB,\quad A>0.
\]
Its original and projected moments are respectively
\begin{equation}\label{r5l:pairmoments}
 2e^{in\theta}\{A\cosh(nx)+iB\sinh(nx)\},
 \qquad 2Ae^{in\theta}.
\end{equation}
Writing the corresponding zero as $\beta+i\gamma$ with
$\beta=1/2-\delta$, one has
\begin{equation}\label{r5l:radialestimates}
 \sinh x=\frac{\delta}{|\rho(1-\rho)|}
 \le\frac{\delta}{\gamma^2+1/4-\delta^2},
 \qquad
 \frac{|B|}{A}=
 \frac{2\delta|\gamma|}{\gamma^2+1/4-\delta^2}
 \le\frac1{|\gamma|}.
\end{equation}
Consequently $x\le\tau_H$ and $|B|/A\le1/H$.  Subtracting
\eqref{r5l:pairmoments} bounds its error by $2A d_n$; this
coefficient is exactly its positive radial mass.  Sum over the
pairs with multiplicities; critical-line terms have zero error.
This proves \eqref{r5l:radialerror}.

The matrix $(b_{k-j})$ is positive semidefinite.  Bound the norm of
its difference from $T_N$ by its maximal absolute row sum to get
\eqref{r5l:radialmatrix}.  For $0\le t\le N\tau_H$,
\[
 \cosh t-1\le\tfrac12t^2\cosh(N\tau_H),
 \qquad \sinh t\le t\cosh(N\tau_H).
\]
Summing $n$ and $n^2$ gives \eqref{r5l:radialmatrixexplicit}.
\end{proof}

\begin{proposition}[The first harmonic changes under radial projection]
\label{r5l:firstmoment}
Under the assumptions of Theorem~\ref{r5l:radialbound},
\[
 a_1\ge b_1,
 \qquad a_1=b_1\quad\Longleftrightarrow\quad\mathrm{RH}.
\]
\end{proposition}
\begin{proof}
Group each off-line quartet with multiplicity $m$, and choose its
member $\rho=\beta+i\gamma$ with $\beta<1/2$ and $\gamma>H$.
Let $u=Re^{i\theta}$, $A=\Re w$, $S=R+R^{-1}$, and
$D=R-R^{-1}$.  Then
\[
 \cos\theta=
 \frac{\gamma^2-\beta(1-\beta)}{|\rho(1-\rho)|}>0,
 \qquad \Im w=-D\sin\theta,
\]
where the second identity follows from $w=2-u-u^{-1}$.
The quartet's exact contribution to $a_1-b_1$ is
\begin{equation}\label{r5l:firstgap}
 2mA(S-2)\cos\theta+2mD^2\sin^2\theta>0.
\end{equation}
Critical-line terms contribute zero, and absolute convergence
permits summation.  Hence the positive quartet contributions cannot
cancel.
\end{proof}

This proposition concerns the specified mass-preserving projection
\eqref{r5l:radialmeasure}.  It does not exclude another positive
measure matching a fixed finite collection of moments.  Matching
the entire original sequence by any positive measure remains
equivalent to RH.

\subsection{Positivity of the original polygons through an explicit degree}

We now obtain actual undamped positivity, rather than an error bound,
on a finite range.  We need the following primary zero-count input.
Bellotti--Wong \cite[Theorem~1.1]{BellottiWong} prove, for $t\ge e$,
\begin{equation}\label{r5l:BW}
 |N(t)-P(t)|\le
 0.11200\log t+0.12567\log\log t+3.77417,
 \qquad P(t)=\frac{t}{2\pi}\log\frac{t}{2\pi e}.
\end{equation}
Here $N(t)$ counts zeros with $0<\Im\rho\le t$, with multiplicity.
In particular,
\begin{equation}\label{r5l:coarsecount}
 |N(t)-P(t)|\le\log t\qquad(t\ge1000).
\end{equation}
Indeed, with $y=\log t$, the difference
$0.888y-0.12567\log y-3.77417$ is positive at $y=\log1000$
and increasing thereafter.

\begin{lemma}[Two populated bands and the full tail mass]
\label{r5l:bands}
For $L\ge2000$,
\begin{equation}\label{r5l:bandcounts}
 N(L)-N(L/2)\ge\frac{L\log L}{8\pi},
 \qquad
 N(4L)-N(2L)\ge\frac{L\log L}{2\pi}.
\end{equation}
For $H\ge20000$,
\begin{equation}\label{r5l:tailmass}
 \sum_{|\Im\rho|>H}|w_\rho|
 \le\sum_{|\gamma|>H}\gamma^{-2}
 <\frac{\log H}{\pi H}.
\end{equation}
\end{lemma}
\begin{proof}
Since $P'(t)=\log(t/(2\pi))/(2\pi)$,
\eqref{r5l:coarsecount} gives
\begin{align*}
 N(L)-N(L/2)&\ge
       \frac{L}{4\pi}\log\frac{L}{4\pi}-2\log L,\\
 N(4L)-N(2L)&\ge
       \frac{L}{\pi}\log\frac{L}{\pi}-\log(8L^2).
\end{align*}
For $L\ge2000$, use
$\log(L/(4\pi))\ge(2/3)\log L$ and
$\log(L/\pi)\ge(3/4)\log L$.  The claimed weaker bounds
follow directly.

For the tail, $|w_\rho|\le\gamma^{-2}$ and partial summation,
including both signs of $\gamma$, gives
\begin{align*}
 \sum_{|\gamma|>H}\gamma^{-2}
 &=-\frac{2N(H)}{H^2}+4\int_H^\infty\frac{N(t)}{t^3}\,dt\\
 &\le\frac{\log(H/(2\pi))+1}{\pi H}
             +\frac{4\log H+1}{H^2}
 <\frac{\log H}{\pi H}.
\end{align*}
The last inequality follows, for example, from
$4\log H+1\le H/5$ and $(\log(2\pi)-1)/\pi>1/5$ in this range.
The factor $2$ in the partial summation is essential.
\end{proof}

\begin{theorem}[A finite range of original polygon positivity]
\label{r5l:finitepositive}
Assume $\mathrm{RH}(H)$ with $H\ge20000$.  For every integer
$1\le M\le H/10$, put $L=\max\{M,2000\}$.  Then, for every
real $\theta$,
\begin{equation}\label{r5l:finitepositivemargin}
 F_M(\theta)\ge\frac{(M+1)\log L}{5000\pi L^3}>0.
\end{equation}
Every original polygon weight consequently satisfies
\begin{equation}\label{r5l:finitepositiveeta}
 \eta_{M,j}\ge
 \frac{(M+1)\log L}{5000\pi L^3(2M+1)}>0.
\end{equation}
\end{theorem}
\begin{proof}
We first give a uniform lower bound for the Fejer polynomial
$F_M^b$ of the positive radial measure.  All zeros in the two
positive-height bands
\[
 \mathcal I_A=(L/2,L],\qquad \mathcal I_B=(2L,4L]
\]
lie on the critical line, since $4L\le H$.  Their weights are
$(1/4+\gamma^2)^{-1}$.  Lemma~\ref{r5l:bands} gives total masses
\begin{equation}\label{r5l:bandmass}
 W_A,W_B\ge w_L:=\frac{\log L}{40\pi L}.
\end{equation}
In fact the respective lower bounds $\log L/(16\pi L)$ and
$\log L/(34\pi L)$ follow by replacing their denominators by
$2L^2$ and $17L^2$.

The Cayley angle of a critical zero at positive height $\gamma$ is
$\phi_\gamma=2\arctan(1/(2\gamma))$.  For
$\gamma_A\in\mathcal I_A$, $\gamma_B\in\mathcal I_B$, put
$a=(M+1)\phi_{\gamma_A}/2$ and
$b=(M+1)\phi_{\gamma_B}/2$.  Then
\begin{equation}\label{r5l:angularseparation}
 \frac{M+1}{5L}\le a-b\le\frac{M+1}{L}<\frac\pi2.
\end{equation}
For the lower bound, use
$\arctan x\ge x-x^3/3$ at $x=1/(2L)$ and
$\arctan x\le x$ at $x=1/(4L)$; the resulting difference is
at least $1/(4L)-1/(24L^3)\ge1/(5L)$.
The upper bound follows from $\arctan(1/L)\le1/L$, together with
$M\le L$ and $L\ge2000$.

For any real $x$ and $d=a-b$ in this range,
\begin{equation}\label{r5l:twophases}
 \sin^2(x-a)+\sin^2(x-b)
 \ge1-\cos d\ge\frac{d^2}{3}
 \ge\frac{(M+1)^2}{75L^2}.
\end{equation}
The middle inequality follows, for instance, from the Taylor
upper bound $\cos d\le1-d^2/2+d^4/24$ and $d<\pi/2<2$.
Take mass $w_L$ from each band measure and couple those equal masses
by their normalized product measure.  This is legitimate even
though individual zero weights are unequal.  The Fejer kernel obeys
\[
 K_M(y)=\frac{\sin^2((M+1)y/2)}{(M+1)\sin^2(y/2)}
 \ge\frac{\sin^2((M+1)y/2)}{M+1},
\]
with the removable points interpreted by continuity.  Integrating
\eqref{r5l:twophases} with $x=(M+1)\theta/2$ and ignoring all
remaining positive mass gives
\begin{equation}\label{r5l:radialfejerlower}
 F_M^b(\theta)\ge
 B_{M,L}:=\frac{(M+1)\log L}{3000\pi L^3}.
\end{equation}

We next control all possible off-line zeros.  Their mass $W$ is at
most $\log H/(\pi H)$ by \eqref{r5l:tailmass}.  For
$n\le M\le H/10$ we have $n\tau_H\le1$.  Therefore
\[
 d_n\le\frac{n^2}{4H^4}+\frac n{H^3},
\]
using $\cosh x-1\le x^2$ and $\sinh x\le2x$ for $0\le x\le1$.
The exact sums
\[
 \sum_{n=1}^M\left(1-\frac n{M+1}\right)n
       =\frac{M(M+2)}6,
 \qquad
 \sum_{n=1}^M\left(1-\frac n{M+1}\right)n^2
       =\frac{M(M+1)(M+2)}{12}
\]
and Theorem~\ref{r5l:radialbound} now give
\begin{align}
 |F_M(\theta)-F_M^b(\theta)|
 &\le W\left\{
 \frac{M(M+1)(M+2)}{24H^4}
       +\frac{M(M+2)}{3H^3}\right\}\nonumber\\
 &\le\frac{\log H\,M(M+1)}{\pi H^4}.
       \label{r5l:fejertailerror}
\end{align}
For the second step, $M\le H/10$ and $M\ge1$ imply
\[
 \frac{(M+1)(M+2)}{24H}+\frac{M+2}{3}\le M+1.
\]

Finally the ratio of \eqref{r5l:fejertailerror} to
$B_{M,L}$ is at most
\begin{equation}\label{r5l:marginratio}
 3000\frac{\log H}{\log L}\frac{ML^3}{H^4}
 \le\frac3{10}\frac{\log H}{\log(H/10)}<\frac25.
\end{equation}
Indeed $M\le L\le H/10$, and $L^4/\log L$ increases for
$L\ge2000$.  The final expression decreases for $H\ge20000$;
at the left endpoint it is less than $2/5$.  Subtracting the
error from \eqref{r5l:radialfejerlower} gives
$F_M>(3/5)B_{M,L}$, proving both conclusions.
\end{proof}

\begin{corollary}[A published height gives an explicit arithmetic range]
\label{r5l:numericrange}
For the actual Riemann zeta function, all the original polygon
weights in \eqref{r5l:fejerdef} are strictly positive for
\[
 1\le M\le300\,000\,000\,000.
\]
\end{corollary}
\begin{proof}
Apply Theorem~\ref{r5l:finitepositive} with the established
Platt--Trudgian height $H=3\cdot10^{12}$.
\end{proof}

This is an analytic consequence of a published finite-height
verification and a published zero-count estimate.  We have not
numerically inspected those polygons.  The full criterion still
requires all $M$ without a bound.  Positivity of these Fejer tests
also does not assert positivity of every arbitrary Toeplitz
quadratic form through the same dimension.  At $\theta=0$ they
include ordinary Li-coefficient positivity, but the assertion here
is uniform in the angle.

\subsection{What finite positive circle averaging cannot remove}

\begin{theorem}[Finite angular averages preserve the radial obstruction]
\label{r5l:angularnogap}
Let $c_1,\ldots,c_J>0$ have sum one, and let
$\theta_1,\ldots,\theta_J$ be fixed angles.  Put
\[
 h_n=\sum_{j=1}^Jc_je^{in\theta_j},
 \qquad \widetilde a_n=h_na_n.
\]
If RH fails and $U=\max_\rho|u_\rho|>1$, then
\begin{equation}\label{r5l:angulargrowth}
 \limsup_{n\to\infty}|\widetilde a_n|^{1/n}=U.
\end{equation}
Consequently such a finite positive angular average cannot make
all Toeplitz matrices positive semidefinite.
\end{theorem}
\begin{proof}
The maximum exists because $|u_\rho|\to1$ with the height, and
the zero symmetries supply a node outside the circle if RH fails.
From \eqref{r5l:Agerm},
\[
 \widetilde A(z)=\sum_jc_jA(e^{i\theta_j}z)
 =\sum_{j,\rho}\frac{c_jw_\rho}{1-u_\rho e^{i\theta_j}z}.
\]
Inside $|z|<1$ the tail is locally normally convergent away from
its poles: $\sum|w_\rho|<\infty$, $u_\rho\to1$, and there are
only finitely many rotations.  At a collision of poles, the
combined coefficient has the form $\sum c_jw_\rho$.  Its real
part is strictly positive by \eqref{r5l:positiveweights}, so no
pole cancels.  At least one pole has modulus $U^{-1}$ and none
has smaller modulus.  Cauchy--Hadamard proves
\eqref{r5l:angulargrowth}.  A positive circle measure of mass
$a_0$ would have all moments bounded by $a_0$, a contradiction.
\end{proof}

For uniform averaging over the $m$th roots of unity this gives
$\limsup_{n\to\infty}|a_{mn}|^{1/n}=U^m$ when RH fails.
Thus a finite cyclic group may delete many frequencies but cannot
delete the radial exponential obstruction.  Finiteness is essential:
uniform averaging over the entire circle deletes every nonconstant
moment.  Continuous Poisson averaging also changes every nonconstant
moment, as its explicit multiplier in Theorem~\ref{r5l:damped} shows.

\subsection{A strict barrier at the nearest pole}

Let
\[
 \tau_*:=\inf\{\tau\ge0:
 (e^{-\tau|k-j|}a_{k-j})_{j,k=0}^N\succeq0
 \text{ for every }N\}.
\]
The allowable set is a closed upper interval: closedness follows
from entrywise limits of each finite matrix, and the upper-interval
property follows by convolution with a Poisson kernel.  It is
nonempty by Theorem~\ref{r5l:damped}.

\begin{proposition}[Convergence damping is not enough]
\label{r5l:strictbarrier}
If RH fails, then
\begin{equation}\label{r5l:strictbarriereq}
 0<\log U<\tau_*\le\tau_H.
\end{equation}
\end{proposition}
\begin{proof}
Damping less than $\log U$ leaves exponential moment growth, so
cannot give a positive circle measure.  At $\tau=\log U$ all
rescaled nodes $v_\rho=u_\rho/U$ have modulus at most one, with
some node $v_0$ on the circle.  If these moments had a positive
representing measure $\nu$, Cesaro atom extraction would give
\begin{equation}\label{r5l:atomextract}
 \nu(\{v_0\})=
 \lim_{K\to\infty}\frac1{K+1}
     \sum_{n=0}^K U^{-n}a_nv_0^{-n}.
\end{equation}
On the zero-sum side, using its scaled-node representation only for
$n\ge0$, dominated convergence by $\sum|w_\rho|$
shows that this limit is $m_0w_{\rho_0}$: the Cayley map is
injective, so only the multiplicity at the chosen node remains.
But $\rho_0$ is off-line and nonreal, and
\[
 \Im w_{\rho_0}=
 \frac{(2\Re\rho_0-1)\Im\rho_0}
             {|\rho_0(1-\rho_0)|^2}\ne0.
\]
It cannot be the mass of a positive-measure atom.  Positivity
therefore fails at $\log U$ as well.  Closedness of the allowable
set gives the strict inequality, and Theorem~\ref{r5l:damped}
gives the upper bound.
\end{proof}

The proposition is a barrier, not a computation of $\tau_*$ for
zeta.  Proposition~\ref{r5l:pairthreshold} computes the exact value
for an individual pair and supplies a concrete model of the gap.
There is also an elementary obstruction to reversing the smoothing
as a universal positivity-preserving Schur operation: its inverse
kernel would be $(r^{-|j-k|})$, whose two-by-two determinant is
$1-r^{-2}<0$ for $r<1$.  An invertible congruence, on the other
hand, preserves inertia rather than repairing it.

\subsection{The remaining arithmetic requirement}

The positive-weight flow can legitimately be applied to the positive
regularized or projected objects established above.  Returning from
them to all original moments requires additional arithmetic control.
The exact first-harmonic loss under the specified radial projection,
the pole persistence under finite angular averages, and the strict
boundary-atom obstruction prove limitations of these particular
operations; they do not prove that every possible method of removing
regularization must fail.

The completed undamped result is the finite interval in
Corollary~\ref{r5l:numericrange}, with the explicit margin
\eqref{r5l:finitepositiveeta}.  The still-open step is a bound
uniformly over all degrees for the actual arithmetic sequence.
No argument here sends the fixed verified height $H$ to infinity.
No historical priority is claimed for the classical moment methods
or for finite-height transfer in general.

\section{The exact global projection onto the nonnegative Nyman cone}
\label{r5g:section}

The local constant $1/92$ from the preceding investigation does not
claim to be the full $L^2(0,1)$ cone distance. We now solve that full
nonnegative-coefficient problem. The result is an explicit finite
projection, certified against all infinitely many generators. Its
role is to close a sign-restricted auxiliary problem; it does not
settle the signed Nyman approximation criterion.

For $n\ge2$, retain
\[
 f_n(x)=\{1/(nx)\}-\{1/x\}/n,\qquad 0<x<1.
\]
Let $\mathcal C_+$ be the cone of all finite sums $\sum_{n\ge2}c_nf_n$
with $c_n\ge0$. We use the real inner product in $L^2(0,1)$.
Define the fixed eight-element set
\begin{equation}\label{r5g:active}
 \mathcal S=\{2,3,4,5,7,11,13,17\}.
\end{equation}
For positive integers $m,n\ge2$, put $\ell=\operatorname{lcm}(m,n)$
and define the exact numbers
\begin{align}
 b_n&=\frac{\log n}{n},\label{r5g:b}\\
 G_{mn}&=\sum_{r=1}^{\ell}
  \frac{(r\bmod m)(r\bmod n)}{mn\ell}
  \left[\psi\left(\frac{r+1}{\ell}\right)
              -\psi\left(\frac r\ell\right)\right],
 \qquad \psi=\Gamma'/\Gamma.\label{r5g:G}
\end{align}
In the rest of this section, $G_{\mathcal S}$ and $b_{\mathcal S}$
mean the eight-by-eight submatrix and corresponding vector in the
increasing order of \eqref{r5g:active}.

\begin{theorem}[Global nonnegative cone optimum]\label{r5g:global}
Set
\begin{equation}\label{r5g:exactopt}
 c^*=G_{\mathcal S}^{-1}b_{\mathcal S},\qquad
 g_* =\sum_{m\in\mathcal S}c_m^*f_m,\qquad
 \mathcal E_+=1-b_{\mathcal S}^{\mathsf T}
                      G_{\mathcal S}^{-1}b_{\mathcal S}.
\end{equation}
Every entry of $c^*$ is strictly positive. The exact global minimum is
\begin{equation}\label{r5g:globalmin}
 \boxed{\displaystyle
 \inf_{g\in\mathcal C_+}\|1-g\|_2^2
 =\min_{g\in\overline{\mathcal C_+}}\|1-g\|_2^2
 =\mathcal E_+.}
\end{equation}
The unique minimizing function is $g_*$. Its only finite
nonnegative-generator representation uses precisely the indices
in \eqref{r5g:active}, with coefficients $c^*$.
A certified enclosure is
\begin{equation}\label{r5g:errorinterval}
 0.0273089278984818<\mathcal E_+<0.0273089278984835.
\end{equation}
The exact answer is the finite expression \eqref{r5g:exactopt};
\eqref{r5g:errorinterval} is a decimal enclosure, not its definition.
\end{theorem}

The proof has an infinite analytic part and a small, rigorous finite
certificate. We give both. The certificate uses only integer arithmetic
and explicit logarithm and digamma remainder bounds; its complete
standard-library Python verifier is reproduced in
Appendix~\ref{r5:cone-code}. No floating-point comparison is used to
decide a sign. No assertion of historical priority is made for the
Hilbert-space projection method or this restricted optimization problem.

\subsection{Exact inner products}

On $I_k=(1/(k+1),1/k)$,
\[
 f_n(x)=v_n(k):=\frac{k\bmod n}{n},\qquad |I_k|=\frac1{k(k+1)}.
\]
Consequently
\begin{equation}\label{r5g:inners}
 \langle1,f_n\rangle=b_n,
 \qquad \langle f_m,f_n\rangle=G_{mn}.
\end{equation}
Here is a direct verification. For the first identity, truncate at
$k=Jn$ and use $k\bmod n=k-n\lfloor k/n\rfloor$ to obtain
\[
 \sum_{k=1}^{Jn}\frac{k\bmod n}{k(k+1)}
 =H_{Jn+1}-H_J-1+\frac{nJ}{Jn+1}\longrightarrow\log n,
\]
where $H_j=\sum_{r=1}^j1/r$. Division by $n$ proves the claim.
For the second, group the absolutely convergent step sum by residues
modulo $\ell$. Its weight at residue $r$ is
\[
 \sum_{j\ge0}\frac1{(j\ell+r)(j\ell+r+1)}
 =\frac1\ell\left[\psi\left(\frac{r+1}{\ell}\right)
                       -\psi\left(\frac r\ell\right)\right],
\]
by the convergent difference of digamma series. This proves
\eqref{r5g:G}.

The generators in any finite set are linearly independent. Indeed,
if $\sum_{n=2}^Nc_nf_n=0$, let $A=\sum c_n/n$ and let $z_k$ be its
step values. Then $z_1=A=0$, and
\[
 z_k-z_{k-1}=A-\sum_{\substack{n\mid k\\n\ge2}}c_n
\]
successively gives $c_k=0$ for $k=2,\ldots,N$. Therefore
$G_{\mathcal S}$ is positive definite and \eqref{r5g:exactopt} is
well defined.

\subsection{A uniform analytic bound for every omitted large index}

For fixed $m\ge2$ define
\[
 \mu_m=\frac{m-1}{2m},\qquad d_m(k)=v_m(k)-\mu_m,
 \qquad
 D_m=\max_{0\le j\le m}\left|\sum_{k=1}^j d_m(k)\right|.
\]
The sequence $d_m$ is periodic with mean zero, so all its partial
sums have absolute value at most $D_m$. Put
\begin{equation}\label{r5g:C}
 C_m=\sum_{k\ge1}\frac{d_m(k)}{k+1}
 =-\frac1m\sum_{r=1}^{m}d_m(r)
                     \psi\left(\frac{r+1}{m}\right).
\end{equation}
The first series converges by summation by parts. Grouping complete
periods and using $\sum_{r=1}^md_m(r)=0$ proves the second identity.

\begin{lemma}[Uniform large-index correlation]\label{r5g:tail}
For every $m,n\ge2$,
\begin{equation}\label{r5g:tailbound}
 \left|nG_{mn}-\mu_m\log n-C_m\right|
 \le\frac{2\zeta(2)D_m}{n}<\frac{10D_m}{3n}.
\end{equation}
\end{lemma}
\begin{proof}
Subtract $\mu_m$ in the step sum and use
$\sum_{k\ge1}(k\bmod n)/[k(k+1)]=\log n$. This gives
\begin{equation}\label{r5g:tailidentity}
 nG_{mn}=\mu_m\log n+C_m
   -n\sum_{k\ge1}d_m(k)e_n(k),
 \qquad e_n(k)=\frac{\lfloor k/n\rfloor}{k(k+1)}.
\end{equation}
Set $e_n(0)=0$. The sequence $e_n$ tends to zero, decreases between
successive multiples of $n$, and jumps upward at $k=jn$.
At each such point,
\[
 0< e_n(jn)-e_n(jn-1)
 =\frac{j(n-2)+1}{jn((jn)^2-1)}
 \le\frac1{jn(jn+1)}\le\frac1{j^2n^2}.
\]
Its total variation is twice the sum of these upward jumps,
and is at most $2\zeta(2)/n^2$.
Summation by parts, using the bounded partial sums of $d_m$,
therefore bounds the last sum in \eqref{r5g:tailidentity} by
$2\zeta(2)D_m/n^2$. This also justifies convergence of the
conditionally convergent pieces used in that identity.
Finally,
\[
 \zeta(2)<\sum_{j=1}^{5}\frac1{j^2}
                       +\int_5^\infty\frac{dt}{t^2}<\frac53.
\]
The displayed strict bound proves the last inequality.
\end{proof}

For any positive coefficient vector on $\mathcal S$, let
$r=1-\sum_{m\in\mathcal S}c_mf_m$ and define
\begin{equation}\label{r5g:abc}
 \alpha=1-\sum_{m\in\mathcal S}c_m\mu_m,
 \quad C_* =\sum_{m\in\mathcal S}c_mC_m,
 \quad D_* =\sum_{m\in\mathcal S}c_mD_m.
\end{equation}
Lemma~\ref{r5g:tail} gives the all-index upper bound
\begin{equation}\label{r5g:dual-tail}
 n\langle r,f_n\rangle
 \le\alpha\log n-C_*+\frac{10D_*}{3n}.
\end{equation}
If $\alpha<0$ and $D_*>0$, the right side is strictly decreasing
as a function of real $n>0$. Thus one negative value at $n=80$
certifies all infinitely many indices $n\ge80$.

\subsection{Finite interval certificate and the global optimality proof}

The integer interval computation uses scale $10^{24}$ and returns
the following outward-rounded bounds for $c^*=G_{\mathcal S}^{-1}
b_{\mathcal S}$:
\begin{center}
\begin{tabular}{@{}rll@{}}
\toprule
Index & Lower bound & Upper bound\\
\midrule
$2$ & $0.87108921237485$ & $0.87108921237486$\\
$3$ & $0.80214904280063$ & $0.80214904280065$\\
$4$ & $0.17794105707088$ & $0.17794105707090$\\
$5$ & $0.53612526970594$ & $0.53612526970596$\\
$7$ & $0.35499581802528$ & $0.35499581802529$\\
$11$ & $0.13208178544897$ & $0.13208178544898$\\
$13$ & $0.04868184161135$ & $0.04868184161136$\\
$17$ & $0.03432498423950$ & $0.03432498423951$\\
\bottomrule
\end{tabular}
\end{center}
All entries are positive. For the residual $r_*=1-g_*$, the
remaining certificate inequalities are
\begin{align}
 \max_{\substack{2\le n\le79\\n\notin\mathcal S}}
           \langle r_*,f_n\rangle&<-0.0002495,
             \label{r5g:finite-signs}\\
 -0.017133038<\alpha&<-0.017133037,\notag\\
 -0.025929730<C_*&<-0.025929729,\notag\\
 1.151255360<D_*&<1.151255361,\notag\\
 \alpha\log80-C_*+10D_*/240&<-0.0011787.
             \label{r5g:tail-sign}
\end{align}
The maximum in \eqref{r5g:finite-signs} occurs at $n=19$.
The exact normal equation gives $\langle r_*,f_m\rangle=0$ for
$m\in\mathcal S$, and \eqref{r5g:dual-tail}--\eqref{r5g:tail-sign}
give strict negativity for every $n\ge80$.
Thus the finite certificate and analytic tail together prove
\begin{equation}\label{r5g:KKT}
 \langle r_*,f_n\rangle
 \begin{cases}=0,&n\in\mathcal S,\\<0,&n\notin\mathcal S.
 \end{cases}
\end{equation}

For clarity, every transcendental error in this certificate is
controlled as follows. For rational $1\le x\le2$, write
$z=(x-1)/(x+1)\in[0,1/3]$. Then
\[
 \log x=2\sum_{j=0}^{34}\frac{z^{2j+1}}{2j+1}+R,
 \qquad 0\le R\le\frac{2z^{71}}{71(1-z^2)}.
\]
Powers of two reduce arbitrary positive rationals to this range.
All operations are rounded outwards to integer multiples of
$10^{-24}$. For digamma values first shift the argument by $16$:
\[
 \psi(x)=\psi(x+16)-\sum_{j=0}^{15}\frac1{x+j}.
\]
For $y=x+16>16$, use the classical real-positive remainder bound
\cite[\S5.11(ii)]{DLMF}:
\begin{equation}\label{r5g:psi-certificate}
 \psi(y)=\log y-\frac1{2y}
           -\sum_{k=1}^{6}\frac{B_{2k}}{2ky^{2k}}+R_\psi,
 \qquad -\frac1{12y^{14}}\le R_\psi\le0.
\end{equation}
This bound also follows by expanding
$(y^2+t^2)^{-1}$ in Binet's positive-kernel integral for $\psi$.
Interval Gaussian elimination has positive pivot lower bounds,
so it encloses the exact vector $c^*$.
The verifier then checks all $70$ excluded indices from $2$ to $79$,
the sign in \eqref{r5g:tail-sign}, and \eqref{r5g:errorinterval}.

\begin{proof}[Completion of the proof of Theorem~\ref{r5g:global}]
For $g=\sum c_nf_n\in\mathcal C_+$, equation \eqref{r5g:KKT}
gives $\langle r_*,g\rangle\le0$ and
$\langle r_*,g_*\rangle=0$. Therefore
\[
 \|1-g\|_2^2
 =\|r_*\|_2^2+\|g-g_*\|_2^2-2\langle r_*,g-g_*\rangle
 \ge\|r_*\|_2^2+\|g-g_*\|_2^2.
\]
The same inequality holds on the norm closure by continuity.
It proves existence and uniqueness of the minimizer and
\eqref{r5g:globalmin}. The normal equation gives
$\|r_*\|_2^2=1-b_{\mathcal S}^{\mathsf T}c^*$.
Finally, equality for a finite nonnegative representation forces
all coefficients outside $\mathcal S$ to vanish by the strict signs
in \eqref{r5g:KKT}; independence of the remaining generators then
forces $c=c^*$. The interval evaluation supplies
\eqref{r5g:errorinterval}.
\end{proof}

This closes the entire nonnegative cone problem, including its
infinite set of possible generators. The signed approximation
space is larger, and the theorem does not give a positive lower
bound for that space. In particular, it does not contradict the
Nyman--Beurling equivalence with RH.

% Round-5 manuscript fragment. No preamble or bibliography.
% Labels: r5n:.
% Bibliography keys: BaezDuarte, Balazard, TaoMobius,
% GallardoSeco2025, TransferZeta, VMRC200.

\section{Cancellation budgets and positive lifts of Nyman flows}
\label{r5n:sec}

We quantify the negative coefficients required by a small signed
Nyman approximation error. The underlying coefficient-recovery
identity is the established Vasyunin--Balazard biorthogonal
construction \cite[Sections~1--2]{Balazard}. Aggregating its finite
functionals gives an explicit obstruction to bounded-mass,
coefficient-preserving positive lifts. We also construct a positive
lift after refining the event set, compute its optimal auxiliary
slope, and exhibit an exact relaxation whose error tends to zero
although none of its members is a finite Nyman combination.
These statements do not establish convergence of the Nyman distances,
and no claim of historical priority is made.

\subsection{The signed approximation and its arithmetic coordinates}
\label{r5n:subsec-setup}

For integers \(n\ge2\), put
\[
 f_n(x)=\left\{\frac1{nx}\right\}
        -\frac1n\left\{\frac1x\right\},
 \qquad 0<x<1.
\]
For a finitely supported real sequence \(c=(c_n)_{n\ge2}\), define
\begin{equation}
\label{r5n:eq-definitions}
 \begin{aligned}
 g_c&=\sum_{n\ge2}c_nf_n,
 &E(c)&=\|1-g_c\|_{L^2(0,1)}^2,\\
 B_0(c)&=\sum_{n\ge2}c_n^-,
 &B_1(c)&=\sum_{n\ge2}\frac{c_n^-}{n},
 \qquad c_n^-=\max\{-c_n,0\}.
 \end{aligned}
\end{equation}
The strong B\'aez-Duarte criterion \cite{BaezDuarte}, with the
tail cancellation recalled in \cite{TransferZeta}, gives
\[
 \mathrm{RH}
 \quad\Longleftrightarrow\quad
 \delta_N:=\inf_{c_2,\ldots,c_N}\sqrt{E(c)}
                   \longrightarrow0.
\]
This is established input. The problem is to prove an unconditional
upper bound tending to zero; the necessary mass bounds below
should not be confused with such an upper bound.

On \(I_k=(1/(k+1),1/k)\), cancellation of the continuous fractional
parts gives
\[
 f_n(x)=\frac{k\bmod n}{n}.
\]
Write
\begin{equation}
\label{r5n:eq-residues}
 z_k=\sum_{n\ge2}c_n\frac{k\bmod n}{n},
 \qquad z_0=0,\qquad w_k=\frac1{k(k+1)}.
\end{equation}
The exact norm is
\begin{equation}
\label{r5n:eq-step-error}
 E(c)=\sum_{k\ge1}w_k(z_k-1)^2.
\end{equation}
Every finite coefficient vector has \(E(c)>0\): if \(L\) is a
common multiple of its support indices, then \(z_L=0\), and the
term \(k=L\) in \eqref{r5n:eq-step-error} is positive.

Let \(\mu\) denote the M\"obius function. The coefficient-recovery
identity \cite[Section~2]{Balazard} is
\begin{equation}
\label{r5n:eq-recovery}
 c_n=-\sum_{d\mid n}\mu(n/d)(z_d-z_{d-1}),
 \qquad n\ge2.
\end{equation}
For completeness, put \(A=\sum c_n/n\). Then
\[
 z_k-z_{k-1}
 =A-\sum_{\substack{n\mid k\\n\ge2}}c_n.
\]
M\"obius inversion proves \eqref{r5n:eq-recovery}, since the
constant \(A\) has zero divisor sum for \(n>1\).
For the formal target \(z_k=1\) when \(k\ge1\), with \(z_0=0\),
the recovered value is \(-\mu(n)\).

\subsection{An aggregate finite cancellation certificate}
\label{r5n:subsec-aggregate}

\begin{proposition}[Aggregate coefficient recovery]
\label{r5n:prop-aggregate}
Fix an integer \(X\ge2\) and nonnegative weights \(a_n\ge0\) for
\(2\le n\le X\). Set
\[
 \begin{aligned}
 T_a&=\sum_{n=2}^Xa_n\mu(n),\\
 b_d&=\sum_{\substack{2\le n\le X\\d\mid n}}a_n\mu(n/d)
            \quad(1\le d\le X),\qquad b_{X+1}=0,\\
 \lambda_k&=b_{k+1}-b_k,\qquad
 K_a=\sum_{k=1}^X k(k+1)\lambda_k^2.
 \end{aligned}
\]
For every finitely supported signed \(c\), without any restriction
on its support relative to \(X\),
\begin{equation}
\label{r5n:eq-aggregate-identity}
 \sum_{n=2}^Xa_nc_n
 =-T_a+\sum_{k=1}^X\lambda_k(z_k-1).
\end{equation}
Consequently,
\begin{equation}
\label{r5n:eq-aggregate-bound}
 \sum_{n=2}^Xa_nc_n^-\ge T_a-\sqrt{K_aE(c)}.
\end{equation}
If \(T_a>B\ge0\) and the left side of
\eqref{r5n:eq-aggregate-bound} is at most \(B\), then
\begin{equation}
\label{r5n:eq-aggregate-error}
 E(c)\ge\frac{(T_a-B)^2}{K_a}.
\end{equation}
\end{proposition}

\begin{proof}
Multiply \eqref{r5n:eq-recovery} by \(a_n\), sum over \(n\le X\),
and collect the terms according to \(d\mid n\):
\[
 \sum_{n=2}^Xa_nc_n
 =-\sum_{d=1}^Xb_d(z_d-z_{d-1})
 =\sum_{k=1}^X(b_{k+1}-b_k)z_k.
\]
The coefficient sum is \(-b_1=-T_a\), proving
\eqref{r5n:eq-aggregate-identity}. Weighted Cauchy--Schwarz and
\eqref{r5n:eq-step-error} give
\[
 \left|\sum_{k=1}^X\lambda_k(z_k-1)\right|
 \le\sqrt{K_aE(c)}.
\]
Since \(-\sum a_nc_n\le\sum a_nc_n^-\), this proves
\eqref{r5n:eq-aggregate-bound}.
Under \(T_a>B\), its rearrangement gives
\eqref{r5n:eq-aggregate-error}; in this case \(K_a>0\), since
\(K_a=0\) would force \(b_1=T_a=0\).
Only finitely many residue cells were used, so coefficients beyond
\(X\) are fully allowed.
\end{proof}

The exact quantity \(K_a\), rather than the upper bounds used below,
is available from a finite divisor calculation. Thus
\eqref{r5n:eq-aggregate-bound} is an explicit finite certificate
for any proposed approximation.

\subsection{Unweighted negative mass}
\label{r5n:subsec-unweighted}

Define
\[
 Q_0(X)=\#\{2\le n\le X:\mu(n)=1\}.
\]
Choose \(a_n=1_{\{\mu(n)=1\}}\) in
Proposition~\ref{r5n:prop-aggregate}. Then \(T_a=Q_0(X)\) and
\[
 |b_d|\le\left\lfloor\frac Xd\right\rfloor\le\frac Xd.
\]
Therefore
\begin{align}
\label{r5n:eq-unweighted-dual}
 K_a
 &\le X^2\sum_{k=1}^X k(k+1)
          \left(\frac1k+\frac1{k+1}\right)^2 \notag\\
 &=X^2\left(4X+1-\frac1{X+1}\right)
 <X^2(4X+1).
\end{align}
Writing \(u_+=\max\{u,0\}\), we obtain the unconditional finite bound
\begin{equation}
\label{r5n:eq-unweighted-finite}
 \boxed{
 E(c)\ge
 \frac{(Q_0(X)-B_0(c))_+^2}{X^2(4X+1)}
 \qquad(X\ge2).}
\end{equation}
No asymptotic theorem is needed for this inequality.

\begin{theorem}[Reciprocal-error negative-mass law]
\label{r5n:thm-unweighted-mass}
For every sequence \(c^{(j)}\) of finitely supported signed
coefficients with \(E_j:=E(c^{(j)})\to0\), one has
\(B_j:=B_0(c^{(j)})\to\infty\), and
\begin{equation}
\label{r5n:eq-unweighted-asymptotic}
 \liminf_{j\to\infty}B_jE_j\ge\frac1{\pi^6}.
\end{equation}
\end{theorem}

\begin{proof}
The elementary squarefree count and the established prime number
theorem give
\begin{equation}
\label{r5n:eq-positive-mobius-count}
 \begin{aligned}
 Q_0(X)
 &=\frac12\sum_{n\le X}\mu(n)^2
    +\frac12\sum_{n\le X}\mu(n)-1\\
 &=\frac3{\pi^2}X+o(X).
 \end{aligned}
\end{equation}
In particular, \(Q_0(X)\to\infty\). If \(B_j\) had a bounded
subsequence, choose a fixed \(X\) for which \(Q_0(X)\) exceeds
that bound. Equation~\eqref{r5n:eq-unweighted-finite} would give
a fixed positive lower bound for the corresponding \(E_j\),
a contradiction. Hence \(B_j\to\infty\).

Take \(X_j=\lfloor\pi^2B_j\rfloor\). By
\eqref{r5n:eq-positive-mobius-count},
\[
 Q_0(X_j)-B_j=2B_j+o(B_j),\qquad
 X_j^2(4X_j+1)=4\pi^6B_j^3(1+o(1)).
\]
Substitution in \eqref{r5n:eq-unweighted-finite} proves
\eqref{r5n:eq-unweighted-asymptotic}.
\end{proof}

\begin{remark}[The unconditional input and the constant]
\label{r5n:rem-pnt}
The squarefree count follows from
\(\mu(n)^2=\sum_{d^2\mid n}\mu(d)\) and
\(\sum\mu(d)/d^2=1/\zeta(2)\). The remaining input in
\eqref{r5n:eq-positive-mobius-count} is
\(\sum_{n\le X}\mu(n)=o(X)\), an established equivalent of PNT,
not an RH assumption. The distinction between PNT and the
elementary bounded harmonic sum used next is recalled in
\cite{TaoMobius}. No rate for PNT is needed here.
The constant \(\pi^{-6}\) is a proved lower bound; its global
optimality is not asserted.
\end{remark}

\subsection{Harmonic negative mass and positive frequency budgets}
\label{r5n:subsec-harmonic}

The weighting \(1/n\) is directly relevant to the flow in
\cite{VMRC200}, because the frequency associated with \(c_n\)
is proportional to \(1/n\). Define
\[
 Q_1(X)=\sum_{\substack{2\le n\le X\\\mu(n)=1}}\frac1n.
\]
Use \(a_n=1_{\{\mu(n)=1\}}/n\) in
Proposition~\ref{r5n:prop-aggregate}. With
\(H_m=\sum_{j=1}^m1/j\),
\begin{equation}
\label{r5n:eq-harmonic-bd}
 |b_d|\le\frac1dH_{\lfloor X/d\rfloor}
       \le\frac{1+\log(X/d)}d.
\end{equation}
Consequently,
\begin{align}
\label{r5n:eq-harmonic-dual}
 K_a
 &\le2\sum_{k=1}^X k(k+1)(b_k^2+b_{k+1}^2)
   =4\sum_{d=1}^Xd^2b_d^2 \notag\\
 &\le4\sum_{d=1}^X(1+\log(X/d))^2 \notag\\
 &\le4\int_0^X(1+\log(X/t))^2\,dt
   =20X.
\end{align}
The equality includes the endpoint \(d=1\), and \(b_{X+1}=0\).
The integral inequality follows because its nonnegative integrand
is decreasing; substitution \(t=Xu\) gives the integral \(5X\).
We have proved the explicit estimate
\begin{equation}
\label{r5n:eq-harmonic-finite}
 \boxed{B_1(c)\ge Q_1(X)-\sqrt{20XE(c)}
       \qquad(X\ge2).}
\end{equation}

\begin{lemma}[Elementary harmonic abundance]
\label{r5n:lem-harmonic-abundance}
One has
\begin{equation}
\label{r5n:eq-harmonic-abundance}
 Q_1(X)=\frac3{\pi^2}\log X+O(1).
\end{equation}
This estimate requires neither PNT nor RH.
\end{lemma}

\begin{proof}
For an integer \(X\ge1\), the divisor identity gives
\[
 1=\sum_{n\le X}\mu(n)\lfloor X/n\rfloor
  =X\sum_{n\le X}\frac{\mu(n)}n
    -\sum_{n\le X}\mu(n)\{X/n\}.
\]
The \(n=1\) fractional part is zero; the remaining at most \(X-1\)
terms have absolute value at most one. Hence
\begin{equation}
\label{r5n:eq-elementary-mobius-harmonic}
 \left|\sum_{n\le X}\frac{\mu(n)}n\right|\le1.
\end{equation}
This elementary inequality is classical; the proof and extensions
are recorded in \cite[Equation~(1.3)]{TaoMobius}.

Also,
\begin{align}
\label{r5n:eq-squarefree-harmonic}
 \sum_{n\le X}\frac{\mu(n)^2}{n}
 &=\sum_{d\le\sqrt X}\frac{\mu(d)}{d^2}
                         H_{\lfloor X/d^2\rfloor}\notag\\
 &=\frac{\log X}{\zeta(2)}+O(1).
\end{align}
Indeed, use
\(H_{\lfloor y\rfloor}=\log y+\gamma+O(1/y)\) uniformly for
\(y\ge1\), together with the absolute convergence of
\(\sum d^{-2}\) and \(\sum(\log d)d^{-2}\).
Truncating the coefficient of \(\log X\) contributes
\(O((\log X)/\sqrt X)\), and the summed harmonic-number error
is \(O(X^{-1/2})\). Thus the bounded remainder requires no PNT.
Finally,
\[
 Q_1(X)
 =\frac12\sum_{n\le X}\frac{\mu(n)^2}{n}
  +\frac12\sum_{n\le X}\frac{\mu(n)}n-1.
\]
Equations~\eqref{r5n:eq-elementary-mobius-harmonic} and
\eqref{r5n:eq-squarefree-harmonic}, with
\(\zeta(2)=\pi^2/6\), prove the result.
\end{proof}

\begin{theorem}[Logarithmic frequency-mass law]
\label{r5n:thm-harmonic-mass}
There is an absolute constant \(C\) such that every finitely
supported signed approximation with \(0<E(c)\le1/2\) satisfies
\begin{equation}
\label{r5n:eq-harmonic-asymptotic}
 B_1(c)\ge\frac3{\pi^2}\log\frac1{E(c)}-C.
\end{equation}
In particular, for every sequence with \(E_j\to0\),
\begin{equation}
\label{r5n:eq-harmonic-liminf}
 \liminf_{j\to\infty}
 \frac{B_1(c^{(j)})}{\log(1/E_j)}\ge\frac3{\pi^2}.
\end{equation}
\end{theorem}

\begin{proof}
Choose \(X=\lfloor1/E(c)\rfloor\). Then \(X\ge2\),
\(XE(c)\le1\), and \(X\ge1/(2E(c))\). By
Lemma~\ref{r5n:lem-harmonic-abundance} and
\eqref{r5n:eq-harmonic-finite},
\[
 \begin{aligned}
 B_1(c)
 &\ge\frac3{\pi^2}\log X-C_0-\sqrt{20}\\
 &\ge\frac3{\pi^2}\log\frac1{E(c)}
       -C_0-\sqrt{20}-\frac3{\pi^2}\log2.
 \end{aligned}
\]
This proves \eqref{r5n:eq-harmonic-asymptotic}.
Dividing by \(\log(1/E_j)\) proves
\eqref{r5n:eq-harmonic-liminf}.
\end{proof}

Equation~\eqref{r5n:eq-harmonic-finite} gives a fully finite
version with no unspecified constant. Its right side uses only
M\"obius values up to the chosen \(X\).
Optimality of the global logarithmic coefficient is not asserted.

\subsection{An exactly sharp local budget curve}
\label{r5n:subsec-local}

The previous note \cite{TransferZeta} determined the local value
\(1/92\) for nonnegative coefficients. The full local
negative-budget curve also has a closed form.

\begin{proposition}[Six-cell budget interpolation]
\label{r5n:prop-local-budget}
For every \(B\ge0\),
\begin{equation}
\label{r5n:eq-local-b0}
 \inf_{\substack{c\ \mathrm{finitely\ supported}\\B_0(c)\le B}}
 \int_{1/7}^1|1-g_c(x)|^2\,dx
 =\frac{(1-B)_+^2}{92}.
\end{equation}
Likewise, for every \(b\ge0\),
\begin{equation}
\label{r5n:eq-local-b1}
 \inf_{\substack{c\ \mathrm{finitely\ supported}\\B_1(c)\le b}}
 \int_{1/7}^1|1-g_c(x)|^2\,dx
 =\frac{(1-6b)_+^2}{92}.
\end{equation}
These are local statements.
\end{proposition}

\begin{proof}
Equation~\eqref{r5n:eq-recovery} at \(n=6\) gives
\[
 c_6=z_3+z_5-z_6-2z_1.
\]
The squared dual norm on the first six weighted cells is
\[
 4\cdot2+12+30+42=92.
\]
Consequently,
\begin{equation}
\label{r5n:eq-six-functional}
 |c_6+1|^2
 \le92\sum_{k=1}^6\frac{(z_k-1)^2}{k(k+1)}.
\end{equation}
If \(B_0(c)\le B\le1\), then \(c_6\ge-B\), proving the lower
bound in \eqref{r5n:eq-local-b0}. If \(B_1(c)\le b\le1/6\),
then \(c_6\ge-6b\), proving the lower bound in
\eqref{r5n:eq-local-b1}.

For sharpness, use the nonnegative endpoint from
\cite{TransferZeta}:
\[
 \begin{aligned}
 c^{(0)}_2&=21/23,&c^{(0)}_3&=19/23,&c^{(0)}_4&=4/23,\\
 c^{(0)}_5&=29/46,&c^{(0)}_6&=0,&c^{(0)}_7&=133/345,
 \end{aligned}
\]
with all other coefficients zero. Direct substitution gives
\[
 (z_1,\ldots,z_6)
 =\left(\frac{22}{23},1,\frac{26}{23},
                1,\frac{61}{46},\frac{25}{46}\right).
\]
Its residual is proportional to the dual vector of
\eqref{r5n:eq-six-functional}, and its local squared error is
\(1/92\).

The second endpoint is the exact six-cell interpolant
\[
 \begin{aligned}
 c^{(1)}_2&=1,&c^{(1)}_3&=1,&c^{(1)}_4&=0,\\
 c^{(1)}_5&=1,&c^{(1)}_6&=-1,&c^{(1)}_7&=14/15,
 \end{aligned}
\]
with all other coefficients zero. Here \(\sum c_n^{(1)}/n=1\).
The divisor-difference identity gives
\(z_1=\cdots=z_6=1\).

For \(0\le t\le1\), take
\(c^{(t)}=(1-t)c^{(0)}+tc^{(1)}\). All its coefficients except
\(c_6=-t\) are nonnegative. Therefore
\[
 B_0(c^{(t)})=t,\qquad B_1(c^{(t)})=t/6.
\]
Its local residual is \(1-t\) times that of \(c^{(0)}\), so its
local squared error is \((1-t)^2/92\).
This proves sharpness in both ranges. For larger budgets,
\(c^{(1)}\) has local error zero.
\end{proof}

\subsection{Coefficient-preserving positive lifts}
\label{r5n:subsec-coordinate-lift}

Let \(L\) be a common multiple of all support indices, and put
\[
 A=\sum_n\frac{c_n}{n},\qquad
 H(s)=\sum_nc_n\{(L/n)s\}-A\{Ls\}.
\]
This is the signed zero-slope flow of \cite{TransferZeta};
the square-jump calculation comes from \cite{VMRC200}.
Its arithmetic step values agree with those of \(g_c\).

If \(c_n=p_n-q_n\) with \(p_n,q_n\ge0\), every such
coefficient-preserving split satisfies
\begin{equation}
\label{r5n:eq-splitting-budget}
 \sum_nq_n\ge B_0(c),\qquad
 \sum_n\frac{q_n}{n}\ge B_1(c).
\end{equation}
Thus Theorems~\ref{r5n:thm-unweighted-mass} and
\ref{r5n:thm-harmonic-mass} constrain every such split, not
only the minimal Jordan decomposition.

For the minimal split, write
\[
 P=\sum_n\frac{c_n^+}{n},\qquad
 Q=\sum_n\frac{c_n^-}{n}=B_1(c),\qquad A=P-Q.
\]
Then \(H=H_+-H_-\), where
\[
 \begin{aligned}
 H_+(s)&=\sum_nc_n^+\{(L/n)s\}+A^-\{Ls\},\\
 H_-(s)&=\sum_nc_n^-\{(L/n)s\}+A^+\{Ls\}.
 \end{aligned}
\]
After deleting zero-weight terms, both coordinates have positive
jump weights and the same continuous slope
\begin{equation}
\label{r5n:eq-positive-slopes}
 C_+=C_-=L\max\{P,Q\}.
\end{equation}
The signed output functional \((u,v)\mapsto u-v\) annihilates
this common slope. Projecting its square identity therefore
returns to the zero-slope identity, not to the positive averaging
step of \cite{VMRC200}.

More generally, a nonconstant zero-slope scalar flow represented
as a linear functional of positive fractional-part flows, each
with positive continuous slope, requires both signs in the
output functional. A nonzero nonnegative functional would have
strictly positive slope. This statement concerns the specified
linear lifts; it does not exclude every nonlinear construction.

The first residue cell has weight \(1/2\) and value \(z_1=A\),
so
\begin{equation}
\label{r5n:eq-first-cell}
 |A-1|\le\sqrt{2E(c)}.
\end{equation}
For \(E(c)<1/2\), one has \(A>0\), hence
\begin{equation}
\label{r5n:eq-normalized-positive-slope}
 \begin{aligned}
 C_+/L=C_-/L
 &=P=A+B_1(c)\\
 &\ge1-\sqrt{2E(c)}
       +\frac3{\pi^2}\log\frac1{E(c)}-C.
 \end{aligned}
\end{equation}
Thus the positive slopes normalized by \(L\) must diverge
logarithmically in this coefficient-preserving lift.
A uniformly bounded auxiliary frequency budget cannot suffice
within this construction. An unbounded-budget lift remains
possible, but its cancellation still requires a separate estimate.

Replacing the signed coefficients by their absolute values does
not restore that estimate. Adding the same positive flow to both
coordinates preserves their difference while increasing their
positive sizes. Moreover, a phase minimizing one positive
coordinate need not minimize the other or the arithmetic weighted
\(L^2\) residual.

\subsection{An exact positive lift after refining the event set}
\label{r5n:subsec-event-lift}

The preceding budgets concern the original indexed generators.
A different positive lift has a different auxiliary cost.

\begin{proposition}[Event-refined positive lift and its optimal slope]
\label{r5n:prop-event-lift}
Let \(H\) be a one-periodic right-continuous real step function with possible jumps
at \(s=k/L\), \(1\le k\le L\). Write
\[
 K_k=H(k/L-)-H(k/L)
\]
for its downward jumps. For any \(\alpha_k\ge|K_k|\), define
\begin{equation}
\label{r5n:eq-event-lift}
 F_\pm(s)=\sum_{k=1}^L
        \frac{\alpha_k\pm K_k}{2}\{s-k/L\}.
\end{equation}
These have nonnegative weights, integral frequency one, and the
same continuous slope
\[
 C_+=C_-=\frac12\sum_{k=1}^L\alpha_k.
\]
They satisfy
\begin{equation}
\label{r5n:eq-event-difference}
 F_+-F_-=H-\int_0^1H(s)\,ds.
\end{equation}
Among all nonnegative two-flow lifts at the fixed events with
these jump differences, the least common slope is exactly
\begin{equation}
\label{r5n:eq-minimal-event-slope}
 C_{\min}=\frac12\sum_k|K_k|.
\end{equation}
\end{proposition}

\begin{proof}
Periodicity gives \(\sum K_k=0\), so the two slopes in
\eqref{r5n:eq-event-lift} coincide. Their difference has the same
jumps and continuous slope as \(H\). Every fractional part in
the sum has mean \(1/2\); hence their difference has mean zero.
This proves \eqref{r5n:eq-event-difference}.
Two nonnegative weights with difference \(K_k\) have sum at least
\(|K_k|\). Summing gives the lower bound in
\eqref{r5n:eq-minimal-event-slope}, attained by
\(\alpha_k=|K_k|\).
Zero-weight terms may be deleted. If strict positivity at every
label is required before deleting terms, the same value is the
infimum, approached by increasing each \(\alpha_k\) slightly.
\end{proof}

For a constant envelope \(\alpha_k=J\ge\max_k|K_k|\), with
\(J>0\), put \(A=F_++F_-\) and center \(H\) as in
\eqref{r5n:eq-event-difference}. Then
\[
 A(s)=J\left(\frac{L-1}{2}+\{Ls\}\right).
\]
The flow \(A\) has slope \(LJ\), jump \(J\), and constant event
midpoint \(a=LJ/2\). If
\(h=(H^-+H^+)/2\) and \(K=H^--H^+\) at an event, direct
expansion gives
\begin{align}
\label{r5n:eq-mixed-cubic}
 \Delta(AH^2)
 &=(a-J/2)(h-K/2)^2-(a+J/2)(h+K/2)^2\notag\\
 &=-2ahK-Jh^2-JK^2/4.
\end{align}
Between events, \((AH^2)'=LJH^2\). Since
\[
 \sum_k h_kK_k
 =\frac12\sum_k\bigl((H(k/L-))^2-(H(k/L))^2\bigr)=0,
\]
summing \eqref{r5n:eq-mixed-cubic} over the period gives exactly
\begin{equation}
\label{r5n:eq-step-quadrature}
 \int_0^1H(s)^2\,ds
 =\frac1L\sum_{k=1}^L
       \left(h_k^2+\frac{K_k^2}{4}\right)
 =\frac1L\sum_{k=1}^LH(k/L)^2.
\end{equation}
Thus the mixed balance is the ordinary step quadrature.
It does not add an upper bound on the particular arithmetic
step values.

There is also a measure distinction.
Equation~\eqref{r5n:eq-step-quadrature} uses uniform phase measure,
whereas the Nyman error uses \eqref{r5n:eq-step-error}.
For an \(L\)-periodic residue sequence, the latter has residue
weights
\[
 W_{L,r}=\sum_{m\ge0}
       \frac1{(mL+r)(mL+r+1)},\qquad 1\le r\le L.
\]
If the step intervals are reparametrized to have these lengths,
the midpoints \(a_k\) of \(A\) vary. The term
\(\sum_k2a_kh_kK_k\) in \eqref{r5n:eq-mixed-cubic} no longer
telescopes away; it retains the missing weighted information.

\subsection{A sharp relaxed countermodel and the finite-support constraint}
\label{r5n:subsec-relaxation}

The next lemma identifies exactly why a convenient positive-flow
relaxation is larger than the set of finite Nyman combinations.

\begin{lemma}[Period and finite index support]
\label{r5n:lem-period-support}
Suppose
\[
 z_k=\sum_{n\ge2}c_n\frac{k\bmod n}{n}
\]
has finite coefficient support and is \(L\)-periodic in the
integer \(k\). Then every index with \(c_n\ne0\) divides \(L\).
\end{lemma}

\begin{proof}
Suppose otherwise, and choose a largest support index \(n\)
not dividing \(L\). Let \(Q\) be a common period of \(L\) and
all support indices. Inspect the discrete Fourier coefficient
at frequency \(1/n\), averaged over \(Q\) terms.
An \(m\)-periodic sequence contributes zero unless \(n\mid m\):
the sum over its repeated blocks is otherwise a vanishing
geometric sum.

A support index \(m>n\) divisible by \(n\) would also fail to
divide \(L\), contrary to the choice of \(n\).
Thus only the \(n\)-th generator can contribute.
Its coefficient is
\[
 \frac1{n^2}\sum_{k=0}^{n-1}k e^{-2\pi i k/n}
 =-\frac1{n(1-e^{-2\pi i/n})}\ne0.
\]
But an \(L\)-periodic sequence has zero coefficient at frequency
\(1/n\) when \(n\nmid L\). This contradicts \(c_n\ne0\).
\end{proof}

\begin{proposition}[Vanishing relaxed error at auxiliary slope one]
\label{r5n:prop-relaxed-countermodel}
For every prime \(p\ge3\), the \(p\)-periodic pattern
\[
 z_k^{(p)}=1_{\{p\nmid k\}},\qquad z_0^{(p)}=0,
\]
has a positive event-refined lift with minimal common slope
exactly one. Its arithmetic weighted error is
\begin{equation}
\label{r5n:eq-relaxed-error}
 E_p=\sum_{k\ge1}
       \frac{|1-z_k^{(p)}|^2}{k(k+1)}
     =\sum_{m\ge1}\frac1{mp(mp+1)}
     \le\frac{\pi^2}{6p^2}.
\end{equation}
None of these patterns is a finite Nyman combination.
Moreover, the least asymptotic auxiliary slope in the unrestricted
periodic residue relaxation with \(z_0=0\) and error tending to
zero is exactly one.
\end{proposition}

\begin{proof}
The pattern has only two nonzero cyclic jumps, one \(+1\) and one
\(-1\). Proposition~\ref{r5n:prop-event-lift} gives \(C_{\min}=1\).
Only multiples of \(p\) contribute to the error, proving the
identity in \eqref{r5n:eq-relaxed-error}; the bound follows from
\(mp(mp+1)\ge m^2p^2\) and \(\sum m^{-2}=\pi^2/6\).

By Lemma~\ref{r5n:lem-period-support}, a finite Nyman representation
of this \(p\)-periodic pattern would use only \(f_p\).
Its values would be proportional to
\[
 1/p,\ 2/p,\ \ldots,\ (p-1)/p,\ 0,
\]
not to \(1,1,\ldots,1,0\).
The values at \(k=1,2\) already give a contradiction for \(p\ge3\).

Finally, any periodic real residue pattern with \(z_0=z_L=0\)
and error \(E\) satisfies \(|z_1-1|\le\sqrt{2E}\).
Its cyclic total variation is at least \(2|z_1|\), so
\begin{equation}
\label{r5n:eq-relaxed-slope-lower}
 C_{\min}\ge|z_1|\ge1-\sqrt{2E}.
\end{equation}
This lower bound and the family \eqref{r5n:eq-relaxed-error}
prove the claimed sharp asymptotic value one.
\end{proof}

Thus an event-refined positive lift can have bounded auxiliary
mass while the relaxed error tends rapidly to zero.
This does not contradict Theorem~\ref{r5n:thm-harmonic-mass}:
the relaxed examples are outside the finite Nyman span.
Their convergence removes the original divisor-compatible
generator constraint. No universal growing-mass obstruction is
being asserted for every conceivable event-refined lift.

\subsection{The arithmetic descent estimate that remains}
\label{r5n:subsec-descent}

We record the standard Hilbert-space mechanism to identify the
remaining estimate precisely; its Schur-complement algebra is
not a new result.
Let \(P_N\) be orthogonal projection onto
\(\operatorname{span}(f_2,\ldots,f_N)\), let \(r_N=1-P_N1\),
and append the block \(N<n\le2N\). Put
\[
 u_n=(I-P_N)f_n,\qquad
 S_{mn}=\langle u_m,u_n\rangle,\qquad
 s_n=\langle r_N,u_n\rangle=\langle r_N,f_n\rangle.
\]
The generators are linearly independent, as follows from
\eqref{r5n:eq-recovery}, so \(S\) is positive definite.
Orthogonal least squares gives
\begin{equation}
\label{r5n:eq-schur-descent}
 \delta_N^2-\delta_{2N}^2=s^{\mathsf T}S^{-1}s.
\end{equation}
For \(e_n(x)=\{1/(nx)\}\) on \((0,\infty)\),
\[
 0\le f_n\le e_n\quad\hbox{on }(0,1),\qquad
 \|e_n\|_{L^2(0,\infty)}^2=\frac{\|e_1\|_2^2}{n}\le\frac2n.
\]
The last bound follows from \(e_1\le1\) on \((0,1)\) and
\(e_1(x)=1/x\) on \((1,\infty)\). Consequently,
\[
 \operatorname{tr}S
 \le\sum_{n=N+1}^{2N}\|f_n\|_2^2
 \le2\sum_{n=N+1}^{2N}\frac1n
 \le2\log2.
\]
Since the largest eigenvalue of \(S\) is at most its trace,
\begin{equation}
\label{r5n:eq-block-descent}
 \delta_N^2-\delta_{2N}^2
 \ge\frac1{2\log2}
       \sum_{n=N+1}^{2N}|\langle r_N,f_n\rangle|^2.
\end{equation}

Suppose one could prove explicit numbers \(\alpha_j\ge0\) such
that
\begin{equation}
\label{r5n:eq-missing-frame}
 \sum_{n=2^j+1}^{2^{j+1}}
       |\langle r_{2^j},f_n\rangle|^2
 \ge\alpha_j\delta_{2^j}^2,
 \qquad \sum_j\alpha_j=\infty.
\end{equation}
Then \eqref{r5n:eq-block-descent} would give
\[
 \delta_{2^J}^2
 \le\delta_{2^{j_0}}^2
       \prod_{j=j_0}^{J-1}
          \left(1-\frac{\alpha_j}{2\log2}\right)
 \longrightarrow0.
\]
The factors lie in \([0,1]\) whenever
\eqref{r5n:eq-missing-frame} holds, by
\eqref{r5n:eq-block-descent}. A zero factor gives the conclusion
immediately; otherwise the divergent sum makes the product zero.
Monotonicity would then imply \(\delta_N\to0\).

These correlations are the arithmetic sums
\begin{equation}
\label{r5n:eq-arithmetic-correlations}
 \langle r_N,f_n\rangle
 =\sum_{k\ge1}\frac{r_N|_{I_k}}{k(k+1)}
                        \frac{k\bmod n}{n}.
\end{equation}
No divergent lower bound in \eqref{r5n:eq-missing-frame} has
been established here. The positive jump identities and the
negative-mass bounds constrain possible successful approximants;
they do not produce the cancellations in
\eqref{r5n:eq-arithmetic-correlations}.

\begin{remark}[Attribution and scope]
\label{r5n:rem-scope}
The fixed coefficient limit \(-\mu(n)\), the step-space framework,
and the biorthogonal recovery mechanism are established results
\cite{Balazard,GallardoSeco2025}. The density and inverse issues
in Noor-type formulations are not reclassified as new results
here. The local nonnegative endpoint \(1/92\) and the signed
zero-slope representation come from \cite{TransferZeta};
Proposition~\ref{r5n:prop-local-budget} supplies its budget
interpolation. The elementary harmonic M\"obius bound is
separate from the stronger PNT input \cite{TaoMobius}.
The conclusions proved here are necessary cancellation budgets,
an exact positive lift and its cost, and a precise relaxed
countermodel. They do not establish an unconditional upper
bound \(\delta_N\le f(N)\to0\).
\end{remark}

\section{The three estimates that still require arithmetic closure}
\label{r5:targets}

The preceding results complete several auxiliary problems in their
full stated domains. Their remaining arithmetic requirements can
be written without an unspecified coordinate change:
\begin{align}
 &\sup_N\sum_{n=2}^Nw_n(\theta_n-\bar\theta_N)^2<\infty,
        \label{r5:target-prime}\\
 &\eta_{M,j}\ge0\quad\hbox{for every }M\ge1,\ 0\le j\le2M,
        \label{r5:target-li}\\
 &\inf_{c_2,\ldots,c_N\in\R}
      \left\|1-\sum_{n=2}^Nc_nf_n\right\|_2\longrightarrow0.
        \label{r5:target-nyman}
\end{align}
Each statement is equivalent to RH by the exact bridges recorded
in Section~\ref{r5:intro}.

For \eqref{r5:target-prime}, every integer lift and common rotation
has been determined, with a uniform separation from alternative
integer lifts. The unbounded corrected cubic upper bound proves
that this particular estimate does not close the target. Its
remaining branch loss must participate quantitatively.

For \eqref{r5:target-li}, the original weights are strictly positive
through degree $3\cdot10^{11}$, and an explicitly damped sequence
is positive at all degrees. Neither assertion removes damping
uniformly over all degrees. The radial first-harmonic discrepancy,
finite-rotation pole persistence, and boundary-atom barrier specify
why the particular inverse operations studied here do not do so.

For \eqref{r5:target-nyman}, the entire nonnegative cone has now
been projected exactly. Signed approximation remains a different
space. The coefficient-preserving lift requires growing negative
budgets, while event refinement permits bounded auxiliary slope
only by allowing patterns outside the original finite span.
An actual divergent lower bound on the residual correlations in
\eqref{r5n:eq-missing-frame} would supply a sufficient descent
mechanism; it is not proved here.

The exact statements proved in this manuscript are independently
checkable mathematical results. They neither certify the full
corpus nor establish a publication-level novelty claim. The two
finite computations are rigorous sign certificates with complete
code, and the infinite tails are handled analytically rather than
discarded by a cutoff.


\appendix
\clearpage

\section{Analytic verification of the original prime-circle criterion}
\label{r4a:appendix}

This appendix verifies the criterion from \cite{PrimeCircles}
independently of the quotient-lattice calculation.  The external
inputs are Mertens' theorem, the prime number theorem, Bertrand's
postulate, the Euler product and functional equation of $\zeta$, and
the following conditional mean-square estimate of Cram\'er
\cite{Cramer}:
\begin{equation}\label{r4a:cramer}
 \mathrm{RH}\quad\Longrightarrow\quad
 \int_X^{2X}(\psi(t)-t)^2\,dt\ll X^2,
 \qquad
 \psi(x)=\sum_{p^r\leq x}\log p.
\end{equation}
The exact form \eqref{r4a:cramer} is stated in
\cite[eq.~(1)]{BrentPlattTrudgian}; their Theorem~1 supplies an
explicit constant.  The pointwise RH estimate for $\psi(x)-x$
alone does not imply \eqref{r4a:cramer}.

All logarithms are natural.  Let
\[
 S(x)=\sum_{p\leq x}\frac1p,
 \qquad E(x)=S(x)-\log\log x-B_{\mathbb P},
 \qquad u_\infty=e^{2\pi iB_{\mathbb P}},
\]
where $B_{\mathbb P}=\lim_{x\to\infty}(S(x)-\log\log x)$.
Mertens' theorem gives $E(x)\to0$.  For $n\geq2$ set
\[
 U_n=e^{2\pi i(S(p_n)-\log\log p_n)},\qquad
 w_n=p_{n+1}-p_{n-1},
\]
and let $W_N=\sum_{n=2}^Nw_n$, $M_N=\sum_{n=2}^Nw_nU_n$,
and $D_N=W_N-|M_N|$.

\begin{theorem}\label{r4a:criterion}
The following statements are equivalent:
\begin{equation}\label{r4a:chain}
 \begin{aligned}
 \mathrm{RH}
 &\quad\Longleftrightarrow\quad E\in L^2([2,\infty),dx)\\
 &\quad\Longleftrightarrow\quad
 \sum_{n=2}^{\infty}w_n|U_n-u_\infty|^2<\infty\\
 &\quad\Longleftrightarrow\quad \sup_ND_N<\infty.
 \end{aligned}
\end{equation}
Under these equivalent conditions,
\begin{equation}\label{r4a:limit}
 2\lim_{N\to\infty}D_N
     =\sum_{n=2}^{\infty}w_n|U_n-u_\infty|^2.
\end{equation}
\end{theorem}

\begin{proof}
\textit{1. RH and the continuous quadratic energy.}
Assume RH, put $\vartheta(x)=\sum_{p\leq x}\log p$, and write
\[
 R(x)=\vartheta(x)-x,\qquad f(x)=\frac{R(x)}{x\log x}.
\]
The elementary estimate
$\psi(x)-\vartheta(x)=O(\sqrt{x})$ follows by summing the
prime-power terms and using $\vartheta(y)\ll y$.
Thus \eqref{r4a:cramer} also holds with $R$ in place of $\psi-t$.
Consequently
\[
 \int_X^{2X}f(x)^2\,dx\ll\frac1{\log^2X}.
\]
Summing over dyadic intervals proves $f\in L^2([2,\infty))$.
Partial summation, with the constant fixed by Mertens' limit, gives
\begin{equation}\label{r4a:hardy}
 E(x)=f(x)-\int_x^\infty
        \frac{f(t)}{t}\left(1+\frac1{\log t}\right)\,dt.
\end{equation}
For completeness, partial summation first gives
\[
 S(x)=\frac{\vartheta(x)}{x\log x}
       +\int_2^x\frac{\vartheta(t)(1+\log t)}{t^2\log^2t}\,dt.
\]
Subtracting its limit after removing $\log\log x$ yields
\eqref{r4a:hardy}.  Here $f(x)\to0$ by the prime number theorem,
and the tail involving $f$ is absolutely convergent by
Cauchy--Schwarz.

The operator $Hg(x)=\int_x^\infty g(t)\,dt/t$ satisfies
$\|Hg\|_2\leq2\|g\|_2$ on $[2,\infty)$.
Indeed, for compactly supported real $g$, integration by parts gives
\[
 \|Hg\|_2^2=-2(Hg(2))^2+2\int_2^\infty(Hg)(x)g(x)\,dx
       \leq2\|Hg\|_2\|g\|_2;
\]
truncation proves the general case.  Since $1+1/\log t$ is bounded
on this interval, \eqref{r4a:hardy} proves $E\in L^2$.

Conversely, suppose $E\in L^2$.  Then
\[
 F(s)=\int_2^\infty E(x)x^{-s}\,dx
\]
is holomorphic for $\Re s>1/2$.  Cauchy--Schwarz gives absolute
convergence uniformly on compact sub-half-planes, also after any
number of differentiations, since
$\int_2^\infty x^{-2\sigma}(\log x)^{2j}\,dx<\infty$ for
$\sigma>1/2$.
For $\Re s>1$ let $P(s)=\sum_p p^{-s}$ and
$c=\log\log2+B_{\mathbb P}$.  Stieltjes integration gives
\[
 P(s)=(s-1)F(s)+c\,2^{1-s}
          +\int_2^\infty\frac{x^{-s}}{\log x}\,dx.
\]
Differentiating and using the Euler product yields
\begin{equation}\label{r4a:logderivative}
 \begin{aligned}
 \frac{\zeta'(s)}{\zeta(s)}
   ={}&F(s)+(s-1)F'(s)-c(\log2)2^{1-s}\\
      &\hspace{1em}-\frac{2^{1-s}}{s-1}
       +\sum_{k\geq2}P'(ks).
 \end{aligned}
\end{equation}
The last sum is holomorphic for $\Re s>1/2$: on every compact
sub-half-plane it is absolutely dominated by
\[
 \sum_p\frac{(\log p)p^{-2\sigma}}{1-p^{-\sigma}}<\infty,
 \qquad \sigma>1/2.
\]
Thus the right side of \eqref{r4a:logderivative} is meromorphic
there with its only possible pole at $s=1$.  By meromorphic
continuation it equals $\zeta'/\zeta$ throughout that half-plane.
A zero of $\zeta$ would produce an additional pole of its
logarithmic derivative, which is impossible.  The functional
equation now gives RH.  No continuation to the boundary
$\Re s=1/2$ is used.

\medskip
\textit{2. Sampling with the actual prime-gap weights.}
Put $h_n=p_{n+1}-p_n$, $e_n=E(p_n)$, and
$b_n=E(p_{n+1}^{-})=e_{n+1}-1/p_{n+1}$.
On $[p_n,p_{n+1})$,
\[
 E(x)=e_n-(\log\log x-\log\log p_n).
\]
The change of variable
$y=(\log\log x-\log\log p_n)/
       (\log\log p_{n+1}-\log\log p_n)$
makes this $(1-y)e_n+yb_n$.  Bertrand's postulate gives a
maximum-to-minimum ratio less than $4$ for $dx/dy$, because
\[
 \frac{p_{n+1}\log p_{n+1}}{p_n\log p_n}<4.
\]
Since $\int_0^1(dx/dy)\,dy=h_n$, we have $dx/dy\geq h_n/4$.
Using
$\int_0^1((1-y)a+yb)^2dy=(a^2+ab+b^2)/3
 \geq(a^2+b^2)/6$
and the pointwise upper bound by $a^2+b^2$, we obtain
\begin{equation}\label{r4a:gap}
 \frac{h_n}{24}(e_n^2+b_n^2)
  \leq\int_{p_n}^{p_{n+1}}E(x)^2\,dx
  \leq h_n(e_n^2+b_n^2).
\end{equation}
The jump correction has finite weighted square:
\[
 \sum_{n\geq1}\frac{h_n}{p_{n+1}^2}
 \leq\int_2^\infty\frac{dx}{x^2}=\frac12.
\]
Hence replacing $b_n$ by $e_{n+1}$ in the sum of
\eqref{r4a:gap} preserves convergence.  Regrouping its endpoint
terms gives weights $h_{n-1}+h_n=w_n$ for $n\geq2$; the initial
endpoint is finite.  Therefore
\[
 E\in L^2\quad\Longleftrightarrow\quad
 \sum_{n\geq2}w_ne_n^2<\infty.
\]
Finally, $e_n\to0$ and
$|U_n-u_\infty|^2=4\sin^2(\pi e_n)\sim4\pi^2e_n^2$,
proving the middle equivalence in \eqref{r4a:chain}.

\medskip
\textit{3. Removing the limiting direction from the statistic.}
Let $J_N=\sum_{n=2}^Nw_n|U_n-u_\infty|^2$.
Weighted Cesaro convergence gives $M_N/W_N\to u_\infty$, since
$U_n\to u_\infty$ and $W_N\to\infty$.
For all sufficiently large $N$, $M_N\ne0$, and its unit direction
$v_N=M_N/|M_N|$ tends to $u_\infty$.  Direct minimization over a
unit vector gives
\begin{equation}\label{r4a:unitminimum}
 2D_N=\min_{|v|=1}\sum_{n=2}^Nw_n|U_n-v|^2
     =\sum_{n=2}^Nw_n|U_n-v_N|^2
     \leq J_N.
\end{equation}
Thus finite $J_\infty$ bounds $D_N$.  Conversely, if $D_N\leq C$,
then for every fixed $K$ and all sufficiently large $N\geq K$,
\[
 \sum_{n=2}^K w_n|U_n-v_N|^2\leq2C.
\]
First let $N\to\infty$ and then $K\to\infty$ to obtain
$J_\infty\leq2C$.  This proves the last equivalence.
The triangle inequality makes $D_N$ nondecreasing.  The same
finite-prefix argument with $2D_N$ on the right gives
$J_\infty\leq2D_\infty$ whenever these quantities are finite;
the reverse inequality follows from \eqref{r4a:unitminimum}.
This proves \eqref{r4a:limit}.
\end{proof}

The theorem validates the analytic criterion used in the main text.
It leaves the assertion $\sup_ND_N<\infty$ unproved, since that
assertion is precisely the equivalent RH target.

\clearpage
\section{Complete prime-arc integer certificate}
\label{r5:prime-code}

The following complete Python 3.9+ program uses only the standard library.
Save the listing as a Python file and execute it normally, without the
optimization flag that disables assertions. A successful run prints a
JSON report recording successful assertions; an assertion failure invalidates
the certificate. Decimal formatting is used only for reporting the
integer bounds. The analytic infinite-tail argument is in the main text.

\begin{lstlisting}
#!/usr/bin/env python3
"""Exact integer-arithmetic certificate for the actual prime phase arc.

No floating-point calculations and no numerical special-function library.
Every bound is an integer multiple of 10^{-18}.  The only infinite-prime
input used after this finite calculation is Rosser--Schoenfeld (1962),
Theorem 5, equations (3.17)--(3.18).
"""

from math import gcd, isqrt
from pathlib import Path
import json

SCALE = 10**18
TERMS = 24
CUTOFF = 10_000


def log_unit_bounds(numerator, denominator):
    """Bounds for log(numerator/denominator), when 1 <= ratio <= 2.

    With z=(ratio-1)/(ratio+1), 0 <= z <= 1/3, use
      log(ratio)=2 sum_{j=0}^{K-1} z^{2j+1}/(2j+1) + R,
      0 <= R <= 9/[4(2K+1)3^{2K+1}] < 1/SCALE.
    Every series term is rounded down using integer division.  The upper
    bound budgets both the 2K rounding units and one tail unit.
    """
    assert 0 < denominator <= numerator <= 2 * denominator
    z_num, z_den = numerator - denominator, numerator + denominator
    common = gcd(z_num, z_den)
    z_num //= common
    z_den //= common
    z_num2, z_den2 = z_num**2, z_den**2
    num_power, den_power = z_num, z_den
    rounded_sum = 0
    for j in range(TERMS):
        rounded_sum += SCALE * num_power // ((2*j + 1) * den_power)
        num_power *= z_num2
        den_power *= z_den2
    assert 9 * SCALE < 4 * (2*TERMS + 1) * 3**(2*TERMS + 1)
    lower = 2 * rounded_sum
    upper = lower + 2 * TERMS + 1
    return lower, upper


LOG2_LO, LOG2_HI = log_unit_bounds(2, 1)


def log_bounds(numerator, denominator=1):
    """Rigorous fixed-point bounds for log of any positive rational."""
    assert numerator > 0 and denominator > 0
    exponent = numerator.bit_length() - denominator.bit_length()
    if exponent >= 0:
        norm_num, norm_den = numerator, denominator << exponent
    else:
        norm_num, norm_den = numerator << (-exponent), denominator
    while norm_num < norm_den:
        norm_num *= 2
        exponent -= 1
    while norm_num >= 2 * norm_den:
        norm_den *= 2
        exponent += 1
    lower, upper = log_unit_bounds(norm_num, norm_den)
    if exponent >= 0:
        return lower + exponent * LOG2_LO, upper + exponent * LOG2_HI
    return lower + exponent * LOG2_HI, upper + exponent * LOG2_LO


def loglog_bounds(integer):
    ln_lo, ln_hi = log_bounds(integer)
    assert ln_lo > 0
    return log_bounds(ln_lo, SCALE)[0], log_bounds(ln_hi, SCALE)[1]


def decimal(integer):
    sign = "-" if integer < 0 else ""
    integer = abs(integer)
    return f"{sign}{integer // SCALE}.{integer % SCALE:018d}"


def main():
    sieve = bytearray(b"\x01") * (CUTOFF + 1)
    sieve[:2] = b"\x00\x00"
    for p in range(2, isqrt(CUTOFF) + 1):
        if sieve[p]:
            sieve[p*p:CUTOFF+1:p] = b"\x00" * ((CUTOFF-p*p)//p + 1)
    primes = [p for p in range(2, CUTOFF+1) if sieve[p]]
    reciprocal_lo = 0
    rows = []
    for count, p in enumerate(primes, 1):
        reciprocal_lo += SCALE // p
        reciprocal_hi = reciprocal_lo + count
        if p == 2:
            continue
        ll_lo, ll_hi = loglog_bounds(p)
        theta_lo = reciprocal_lo - ll_hi
        theta_hi = reciprocal_hi - ll_lo
        assert theta_lo * 4 > SCALE
        assert theta_hi * 50 < 37 * SCALE
        rows.append((p, theta_lo, theta_hi))

    llT_lo, llT_hi = loglog_bounds(CUTOFF)
    thetaT_lo = reciprocal_lo - llT_hi
    thetaT_hi = reciprocal_lo + len(primes) - llT_lo
    logT_lo, _ = log_bounds(CUTOFF)
    # delta >= 1/log(T)^2 = 2 epsilon_T.
    denominator = logT_lo**2
    delta = (SCALE**3 + denominator - 1) // denominator
    tail_lo, tail_hi = thetaT_lo - delta, thetaT_hi + delta
    assert tail_lo * 4 > SCALE
    assert tail_hi * 50 < 37 * SCALE

    min_row = min(rows, key=lambda row: row[1])
    max_row = max(rows, key=lambda row: row[2])
    assert min_row[0] == 8597
    assert max_row[0] == 3
    # Also certify that the indicated extrema really are unique.
    assert all(row[1] > min_row[2] for row in rows if row[0] != min_row[0])
    assert all(row[2] < max_row[1] for row in rows if row[0] != max_row[0])

    result = {
        "method": "exact integer arithmetic; atanh logarithm series with rigorous tail",
        "fixed_point_denominator": SCALE,
        "series_terms": TERMS,
        "prime_cutoff": CUTOFF,
        "prime_count_including_2": len(primes),
        "finite_phase_count": len(rows),
        "finite_minimum": {"prime": min_row[0], "lower": decimal(min_row[1]), "upper": decimal(min_row[2])},
        "finite_maximum": {"prime": max_row[0], "lower": decimal(max_row[1]), "upper": decimal(max_row[2])},
        "theta_at_cutoff": [decimal(thetaT_lo), decimal(thetaT_hi)],
        "tail_phase_enclosure_RS": [decimal(tail_lo), decimal(tail_hi)],
        "all_prime_phase_interval": ["1/4", "37/50"],
        "all_prime_phase_width_strict_upper": "49/100",
        "certificate_passed": True,
    }
    output = Path(__file__).with_suffix(".json")
    output.write_text(json.dumps(result, indent=2) + "\n")
    print(json.dumps(result, indent=2))


if __name__ == "__main__":
    main()
\end{lstlisting}

\clearpage
\section{Complete global cone integer-interval certificate}
\label{r5:cone-code}

The following complete Python 3.9+ program uses only the standard library.
Save the listing as a Python file and execute it normally, without the
optimization flag that disables assertions. A successful run prints a
JSON report recording successful assertions; an assertion failure invalidates
the certificate. Decimal formatting is used only for reporting the
integer bounds. The analytic infinite-tail argument is in the main text.

\begin{lstlisting}
#!/usr/bin/env python3
"""Integer-interval certificate for a global nonnegative Nyman cone projection.
Only Python's standard library is used. No floating-point decision occurs.
"""
from functools import lru_cache
from fractions import Fraction
from math import gcd, lcm
from decimal import Decimal, localcontext
import json
import time

SCALE = 10**24
class I:
    __slots__ = ('lo','hi')
    def __init__(self, lo, hi=None):
        self.lo=int(lo); self.hi=int(lo if hi is None else hi)
        assert self.lo <= self.hi
    @classmethod
    def rat(cls,p,q=1):
        assert q != 0
        return cls(p*SCALE//q, -((-p*SCALE)//q))
    def __add__(self,b):
        if not isinstance(b,I): b=I.rat(b)
        return I(self.lo+b.lo,self.hi+b.hi)
    __radd__=__add__
    def __neg__(self): return I(-self.hi,-self.lo)
    def __sub__(self,b): return self+(-b if isinstance(b,I) else -I.rat(b))
    def __rsub__(self,b): return I.rat(b)+(-self)
    def __mul__(self,b):
        if not isinstance(b,I): b=I.rat(b)
        v=[x*y for x in (self.lo,self.hi) for y in (b.lo,b.hi)]
        return I(min(v)//SCALE,-((-max(v))//SCALE))
    __rmul__=__mul__
    def __truediv__(self,b):
        if not isinstance(b,I): b=I.rat(b)
        assert not (b.lo <= 0 <= b.hi)
        lows=[x*SCALE//y for x in (self.lo,self.hi) for y in (b.lo,b.hi)]
        highs=[-((-x*SCALE)//y) for x in (self.lo,self.hi) for y in (b.lo,b.hi)]
        return I(min(lows),max(highs))
    def report(self):
        with localcontext() as ctx:
            ctx.prec=40
            return [str(Decimal(self.lo)/Decimal(SCALE)),str(Decimal(self.hi)/Decimal(SCALE))]

def log_unit(p,q):
    """log(p/q), 1 <= p/q <= 2, using atanh and a positive tail bound."""
    assert q <= p <= 2*q
    z=I.rat(p-q,p+q); z2=z*z; power=z; value=I.rat(0)
    terms=35
    for j in range(terms):
        value=value+power/(2*j+1)
        power=power*z2
    tail=(2*power)/((2*terms+1)*(1-z2))
    return 2*value+I(0,tail.hi)
LOG2=log_unit(2,1)
@lru_cache(None)
def lograt(p,q=1):
    assert p>0 and q>0
    d=gcd(p,q);p//=d;q//=d
    k=0
    while p>=2*q: q*=2;k+=1
    while p<q: p*=2;k-=1
    return k*LOG2+log_unit(p,q)

BERNOULLI=((1,6),(-1,30),(1,42),(-1,30),(5,66),(-691,2730))
@lru_cache(None)
def psi_reduced(p,q):
    """Digamma at p/q; six-term Binet expansion with signed remainder."""
    assert p>0 and q>0
    shift=16
    xp=p+shift*q
    inv=I.rat(q,xp); inv2=inv*inv
    value=lograt(xp,q)-inv/2
    power=inv2
    for j,(bn,bd) in enumerate(BERNOULLI,1):
        value=value-I.rat(bn,bd*(2*j))*power
        power=power*inv2
    # B_14/(14*x^14) = 1/(12*x^14); the omitted first term is negative.
    rem=power/12
    value=value+I(-rem.hi,0)
    for j in range(shift): value=value-I.rat(q,p+j*q)
    return value

def psi(p,q):
    d=gcd(p,q)
    return psi_reduced(p//d,q//d)

@lru_cache(None)
def gram(m,n):
    if m>n: return gram(n,m)
    period=lcm(m,n); value=I.rat(0)
    for r in range(1,period+1):
        num=(r%m)*(r%n)
        if num:
            weight=(psi(r+1,period)-psi(r,period))/period
            value=value+I.rat(num,m*n)*weight
    return value

def solve(A,b):
    A=[row[:]+[v] for row,v in zip(A,b)];n=len(b)
    for k in range(n):
        pivot=A[k][k]
        assert pivot.lo>0
        for j in range(k+1,n+1): A[k][j]=A[k][j]/pivot
        A[k][k]=I.rat(1)
        for i in range(n):
            if i==k: continue
            factor=A[i][k]
            for j in range(k+1,n+1): A[i][j]=A[i][j]-factor*A[k][j]
            A[i][k]=I.rat(0)
    return [row[-1] for row in A]

def run():
    start=time.monotonic()
    active=(2,3,4,5,7,11,13,17)
    G=[[gram(m,n) for n in active] for m in active]
    b=[lograt(n)/n for n in active]
    c=solve(G,b)
    assert all(v.lo>0 for v in c)
    margins={}
    for n in range(2,80):
        if n not in active:
            v=lograt(n)/n-sum((cm*gram(m,n) for m,cm in zip(active,c)),I.rat(0))
            assert v.hi<0,(n,v.report())
            margins[n]=v
    C=[];D=[];means=[]
    for m in active:
        mean=Fraction(m-1,2*m);means.append(I.rat(mean.numerator,mean.denominator))
        ck=I.rat(0);partial=Fraction(0);maxpartial=Fraction(0)
        for r in range(1,m+1):
            dk=Fraction(r%m,m)-mean
            ck=ck-I.rat(dk.numerator,dk.denominator)*psi(r+1,m)/m
            partial+=dk;maxpartial=max(maxpartial,abs(partial))
        assert partial==0
        C.append(ck);D.append(I.rat(maxpartial.numerator,maxpartial.denominator))
    alpha=1-sum((v*w for v,w in zip(c,means)),I.rat(0))
    sc=sum((v*w for v,w in zip(c,C)),I.rat(0))
    sd=sum((v*w for v,w in zip(c,D)),I.rat(0))
    tail80=alpha*lograt(80)-sc+I.rat(10,3)*sd/80
    assert alpha.hi<0 and sd.lo>0 and tail80.hi<0
    error=1-sum((v*w for v,w in zip(c,b)),I.rat(0))
    assert error.lo>0
    worst=max(margins,key=lambda n:margins[n].hi)
    assert all(margins[worst].lo>v.hi for n,v in margins.items() if n!=worst)
    assert all(v.hi<I.rat(-2495,10**7).lo for v in margins.values())
    assert tail80.hi<I.rat(-11787,10**7).lo
    assert error.lo>I.rat(273089278984818,10**16).hi
    assert error.hi<I.rat(273089278984835,10**16).lo
    result={
      'status':'PASS','arithmetic':'outward-rounded integer intervals, scale 10^24',
      'active':active,'coefficient_intervals':{str(n):v.report() for n,v in zip(active,c)},
      'squared_distance_interval':error.report(),
      'inactive_range':[2,79], 'inactive_count':len(margins),
      'largest_inactive_gradient_index':worst,
      'largest_inactive_gradient_interval':margins[worst].report(),
      'all_inactive_gradient_intervals':{str(n):v.report() for n,v in margins.items()},
      'alpha_interval':alpha.report(),'weighted_C_interval':sc.report(),
      'weighted_D_interval':sd.report(),'tail_bound_at_80_interval':tail80.report(),
      'log_algorithm':'35 atanh terms plus exact geometric tail upper bound',
      'psi_algorithm':'shift 16; B_2 through B_12; negative B_14 remainder bound',
      'elapsed_seconds':round(time.monotonic()-start,3),
    }
    print(json.dumps(result,indent=2))

if __name__=='__main__': run()
\end{lstlisting}

\clearpage
\addcontentsline{toc}{section}{References}
\begin{thebibliography}{99}

\bibitem{VMRC200}
Shaitetsu,
\emph{Weighted fractional-part flows and sharp lattice covering bounds},
Verifiable Mathematical Reasoning Corpus, manuscript [200],
18 November 2026. Checked repository revision
\texttt{ac1492a50784be919fe0b5e8f19f96f8d221a22e}.
\url{https://github.com/QEEDDD/Verifiable-Mathematical-Reasoning-Corpus}.


\bibitem{PrimeCircles}
\emph{Prime Reciprocal Circles and the Riemann Zeta Function:
A Unified Energy, Operator, and Sign Theory},
research synthesis prepared for Keio,
Verifiable Mathematical Reasoning Corpus, manuscript [149],
9 November 2026. Same checked repository revision as [200].
\url{https://github.com/QEEDDD/Verifiable-Mathematical-Reasoning-Corpus}.


\bibitem{BaezDuarte}
L.~B\'aez-Duarte,
A strengthening of the Nyman--Beurling criterion for the Riemann hypothesis,
\emph{Rend. Lincei Mat. Appl.} (9) \textbf{14} (2003), no.~1, 5--11.
Checked preprint: arXiv:math/0202141v2.
\url{https://arxiv.org/abs/math/0202141}.


\bibitem{Balazard}
M.~Balazard,
\emph{An arithmetical function related to B\'aez-Duarte's criterion
for the Riemann hypothesis}, arXiv:1812.04309v1 (2018),
especially Sections~1--2 and the attribution there to V.~I.~Vasyunin.
\url{https://arxiv.org/abs/1812.04309}.


\bibitem{Lagarias}
J.~C.~Lagarias,
On a positivity property of the Riemann $\xi$-function,
\emph{Acta Arith.} \textbf{89} (1999), no.~3, 217--234,
\url{https://doi.org/10.4064/aa-89-3-217-234};
with correction, \emph{Acta Arith.} \textbf{116} (2005), 293--294.
The positive-real criterion is stated in the original introduction
and credited there to earlier work of A.~Hinkkanen.
Author version:
\url{https://websites.umich.edu/~lagarias/doc/positivity.pdf}.
The results here do not use the corrected Lemma~3.1.


\bibitem{BrentPlattTrudgian}
R.~P.~Brent, D.~J.~Platt, and T.~S.~Trudgian,
The mean square of the error term in the prime number theorem,
\emph{J. Number Theory} \textbf{238} (2022), 740--762.
\url{https://doi.org/10.1016/j.jnt.2021.09.016}.
The qualitative conditional input used here is equation~(1), p.~741.
Author-hosted published version:
\url{https://maths-people.anu.edu.au/~brent/pd/rpb274-final.pdf}.


\bibitem{Cramer}
H.~Cram\'er,
Ein Mittelwertsatz in der Primzahltheorie,
\emph{Math. Z.} \textbf{12} (1922), 147--153.
\url{https://doi.org/10.1007/BF01482072}.
The precise dyadic form used here was checked in
Brent--Platt--Trudgian, cited above.


\bibitem{TransferZeta}
Shaitetsu,
\emph{Lattice Covering Flows and the Riemann Zeta Function:
Exact Quadratic Coverings, Polygon Positivity, and Signed Approximation},
research manuscript prepared for Keio, 20 November 2026.
Preceding 28-page research batch, source file
\texttt{vmrc200\_zeta\_circle\_transfers.tex}.
Its principal criteria and normalizations are restated here.

\bibitem{RosserSchoenfeld}
J.~B. Rosser and L. Schoenfeld,
Approximate formulas for some functions of prime numbers,
\emph{Illinois J. Math.} \textbf{6} (1962), no.~1, 64--94.
\url{https://doi.org/10.1215/ijm/1255631807}.
Theorem~5, equations~(3.17)--(3.18), p.~70.
Checked original scan:
\url{https://denisevellachemla.eu/Rosser-Schoenfeld-1962.pdf}.

\bibitem{PlattTrudgian}
D.~Platt and T.~Trudgian,
The Riemann hypothesis is true up to $3\cdot10^{12}$,
\emph{Bull. London Math. Soc.} \textbf{53} (2021), no.~3, 792--797.
\url{https://doi.org/10.1112/blms.12460}.
Checked author preprint: \url{https://arxiv.org/abs/2004.09765}.

\bibitem{BellottiWong}
C.~Bellotti and P.-J.~Wong, with an appendix by A.~Fiori,
Improved estimates for the argument and zero-counting function
of the Riemann zeta-function,
\emph{Math. Comp.}, published electronically 14 November 2025.
\url{https://doi.org/10.1090/mcom/4133}.
Checked author version: arXiv:2412.15470v2, 7 July 2025,
Theorem~1.1, second estimate.
\url{https://arxiv.org/abs/2412.15470}.

\bibitem{TaoMobius}
T.~Tao,
A remark on partial sums involving the M\"obius function,
\emph{Bull. Aust. Math. Soc.} \textbf{81} (2010), 343--349.
\url{https://doi.org/10.1017/S0004972709000884}.
Equation~(1.3) recalls the classical elementary harmonic bound;
the stronger PNT equivalence is distinguished in the introduction.

\bibitem{GallardoSeco2025}
E.~A. Gallardo-Guti\'errez and D.~Seco,
Zero-Free Regions of the Riemann Zeta Function and Approximation
in Weighted Dirichlet Spaces,
\emph{Complex Anal. Oper. Theory} \textbf{19} (2025), Article~38.
\url{https://doi.org/10.1007/s11785-025-01661-2}.
Referenced for context and prior coefficient-limit formulations;
its generator normalization differs from the one used here.

\bibitem{DLMF}
NIST Digital Library of Mathematical Functions,
\S5.9, equation~5.9.15, and \S5.11(ii), error bounds for the
real-positive digamma asymptotic expansion.
\url{https://dlmf.nist.gov/5.9};
\url{https://dlmf.nist.gov/5.11}.
The sign and magnitude of the first omitted term are used in
the integer-interval certificate.


\end{thebibliography}
\end{document}
